\chapter{Vertex algebras}
\section{Vector superspaces}
\begin{dang}
    \begin{definition}\label{Definition: definition of vector superspaces}\cite{de2006finite}
        A vector superspace $V$ is a $\bb{Z}/(2)-$graded vector space. We write $V=V_{\Bar{0}}\oplus V_{\Bar{1}}$ and if $v\in V_\alpha$ is a homogeneous vector of $V$, we define its \textit{parity} by $p(v)=\alpha$. A subspace $W$ of a vector superspace $V$ is a subspace respecting the $\bb
        Z/(2)-$gradation of $V$, i.e. $W=W_{\Bar{0}}\oplus W_{\Bar{1}}$ with $W_\alpha\subset V_\alpha\;\;\forall\alpha\in\bb{Z}/(2)$.
    \end{definition}
\end{dang}
It's important to note that if $V, W$ are two vector superspaces, then $V\otimes W$ can be given a natural structure of a vector superspace by defining its $\bb{Z}/(2)-$gradation as follows:
\begin{equation}\label{Equation: grading in the tensor product of superspaces}
    (V\otimes W)_\alpha=\bigoplus_{p+q=\alpha}V_p\otimes W_q.
\end{equation}
So for homogeneous vectors $v\in V,w\in W$ we have $p(v\otimes w)=p(v)+p(w)$.
\begin{dang}
    \begin{definition}\cite{kac2013bombay}
        An associative superalgebra $A$ is a $\bb{Z}/(2)-$graded associative algebra with the additional property that for $a\in A_{\alpha},\;b\in A_\beta\Rightarrow ab\in A_{\alpha+\beta}$.
    \end{definition}
\end{dang}

\begin{dang}
    \begin{definition}\cite{kac2013bombay}
        A Lie superalgebra $\mathfrak{g}=\mathfrak{g}_{\Bar{0}}\oplus\mathfrak{g}_{\Bar{1}}$ is a $\bb{Z}/(2)-$graded vector superspace equipped with a bilinear map $[\;,\;]: \mathfrak{g}\times \mathfrak{g}\rightarrow \mathfrak{g}$ such that the following properties are satisfied for homogeneous elements $x,y\in \mathfrak{g}$:
        \begin{enumerate}
            \item (super-skewsymmetry) $[y,x]=-p(x,y)[x,y]$ 
            \item (super-Jacobi identity) $[x,[y,z]]=[[x,y],z]+p(x,y)[y,[x,z]]$ 
        \end{enumerate}
        where $p(x,y)=(-1)^{p(x)p(y)}$. We also impose the condition $p([x,y])=p(x)+p(y)$.
    \end{definition}
\end{dang}
As a consequence of the definition, we see that the restriction $[\;,\;]: \mathfrak{g}_{\Bar{1}}\times \mathfrak{g}_{\Bar{1}}\rightarrow \mathfrak{g}_{\Bar{0}}$ is symmetric while $[\;,\;]: \mathfrak{g}_{\Bar{0}}\times \mathfrak{g}_{\Bar{0}}\rightarrow \mathfrak{g}_{\Bar{0}}$ is skewsymmetric. Indeed $\mathfrak{g}_0$ is a Lie algebra in the usual sense. For Lie superalgebras we have an analogous definition for the universal enveloping algebra
\begin{equation}
    \mathcal{U}(\mathfrak{g})=\frac{T(\mathfrak{g})}{\langle x\otimes y-p(x,y)y\otimes x-[x,y]\rangle}
\end{equation}
where $T(\mathfrak{g})=\bigoplus_{k\geq0}T^k(\mathfrak{g})$ is the tensor algebra of $\mathfrak{g}$. It's important to note that $T(\mathfrak{g})$ is a vector superspace according to the grading~(\ref{Equation: grading in the tensor product of superspaces}) and that the ideal $\langle  x\otimes y-p(x,y)y\otimes x-[x,y]\rangle$ is generated by homogeneous elements, so it's a subspace of the vector superspace $T(\mathfrak{g})$. Therefore, the quotient $\mathcal{U}(\mathfrak{g})$ is also a vector superspace with grading 
\begin{equation}
    p(x_1\cdots x_k)=p(x_1)+\cdots+p(x_k)\quad x_i\in\mathfrak{g}.
\end{equation}
In particular this proves that $\mathcal{U}(\mathfrak{g})$ is an associative superalgebra.\par
There is an analogous version of the PBW-theorem for Lie superalgebras.
\begin{dang}
    \begin{theorem}\label{Theorem: PBW theorem for Lie superalgebras}\cite{kac2013bombay}
        Let $\mathfrak{g}=\mathfrak{g}_{\Bar{0}}\oplus\mathfrak{g}_{\Bar{1}}$ be a Lie superalgebra, $\{x_1,\dots,x_p\}$ be a basis for $\mathfrak{g}_{\Bar{0}}$ and let $\{y_1,\dots,y_q\}$ be a basis for $\mathfrak{g}_{\Bar{1}}$. Then $\mathcal{U}(\mathfrak{g})$ is spanned by 
        \begin{equation}
        x_1^{r_1}\cdots x_p^{r_p} y_1^{s_1}\cdots y_q^{s_q}\quad r_i\geq0,\;s_i\leq 1 .     
        \end{equation}
    \end{theorem}
\end{dang}
Given an associative superalgebra $A$, one can endow it with a Lie superalgebra structure by defining the bracket $[a,b]=ab-p(a,b)ba$. As an example, let $V=V_{\Bar{0}}\oplus V_{\Bar{1}}$ be a vector superspace. Then $\text{End}V$ has an associative superalgebra structure by defining the grading $\text{End}V=(\text{End}(V))_{\Bar{0}}\oplus(\text{End}(V))_{\Bar{1}}$ where $(\text{End}(V))_{\alpha}=\{T\in \text{End}(V)\;| TV_{\beta}\subset V_{\alpha+\beta}\}$. This can be seen by choosing a basis for $V$, so we can identify 
\begin{equation}
    (\text{End}(V))_{\Bar{0}}=\left\{
    \begin{pmatrix}
        *&0\\0&*
    \end{pmatrix}\right\}\;\; \text{ and }\;\;
     (\text{End}(V))_{\Bar{1}}=\left\{
    \begin{pmatrix}
        0&*\\ *&0
    \end{pmatrix}\right\}.
\end{equation}


\section{The first definition}
\begin{dang}
    \begin{definition}
        Let $U$ be a vector superspace space, a $U-$valued formal distribution in $l$ variables is a formal series $a(z_1,\dots,z_l)=\sum_{n\in\bb{Z}}a_{n_1,\dots,n_l}z_1^{-n_1-1}\cdots z_l^{-n_l-1}$ with coefficients $a_{n_1,\dots,n_l}\in U$. We use the notation $U[[z_1^{\pm},\dots,z_l^{\pm}]]$ for the \textit{vector superspace space} of $U-$valued formal distributions in $l$ variables.
    \end{definition}
\end{dang}

\begin{dang}
    \begin{definition}\cite{de2006finite}
        Given a vector superspace $V$, a \textit{quantum field} on $V$ is an $\text{End}V-$valued formal distribution $a(z)=\sum_{n\in\bb
        Z}a_{(n)}z^{-n-1}$ with the property that for each $v\in V$ there is a positive integer $N=N(v)\geq0$ such that $a_{(n)}v=0\;\forall n\geq N$.
    \end{definition}
\end{dang}
\begin{dang}
    \begin{definition}\cite{de2006finite}
        Let $V$ be a vector superspace. We say that a pair $(a(z),b(z))$ of $\text{End}V-$valued quantum field is \textit{local} if 
        \begin{equation}
            (z-w)^N[a(z),b(w)]=0
        \end{equation}
        for some $N\geq 0$.
    \end{definition}
\end{dang}
\begin{dang}
    \begin{definition}
        A vertex algebra is a quadruple $(V,T, |0\rangle,\mathcal{F})$ where $V$ is a vector superspace, $T$ is an endomorphism of $V$ called the \textit{infinitesimal translation operator} (or just translation operator), $|0\rangle$ is an element of $V_{\Bar{0}}$ called \textit{vacuum} and $\mathcal{F}=\{a^j(z)\}_{j\in J}$ is a collection of $\text{End}V-$valued quantum fields with the property that all endomorphisms $a_n^j$ have the same parity for a fixed $j\in J$. This data satisfies the following axioms:
        \begin{enumerate}
            \item[(1)] (Vacuum and translation covariance) $T|0\rangle = 0$ and $[T,a^j(z)]=\partial_za^j(z)\;\forall j\in J$.
            \item[(2)] (Locality) $(z-w)^{N_{jk}}[a^j(z),a^k(w)]=0$ for some $N_{jk}\geq0$.
            \item[(3)] (Completeness) $V$ is spanned by element of the form $a_{(n_1)}^{j_1}\cdots a_{(n_k)}^{j_k}|0\rangle$
        \end{enumerate}
    \end{definition}
\end{dang}
It's important to note that the translation operator is uniquely defined by axioms (1) and (3). In fact, by noting that $[T,a(z)]=\partial_za(z)\iff [T,a_{(n)}]=-na_{(n-1)}$, we get
\begin{equation}\label{Equation: translation operator on V}
    Ta_{n_1}^{j_1}\cdots a_{n_k}^{j_k}|0\rangle=\sum_{i=1}^k-n_ia_{n_1}^{j_1}\cdots a_{(n_{i-1})}^{j_{i-1}}a_{(n_i-1)}^{j_i}a_{(n_{i+1})}^{j_{i+1}}\cdots  a_{n_k}^{j_k}|0\rangle
\end{equation}
because of the axiom (1). Now, completeness (axiom (3)) implies that $T$ is fully determined in $V$ by the formula~(\ref{Equation: translation operator on V}).
\begin{dang}
    \begin{lemma}\label{Lemma: a(z)|0> only have positive powers of z}\cite{de2006finite}
        Let $V$ be a vector superspace and let $|0\rangle \in V_{\Bar{0}}, T\in(\text{End}V)_{\Bar{0}}$ be such that $T|0\rangle =0$. Then for any $T-$covariant $\text{End}V-$valued quantum field $a(z)$, we have $a(z)|0\rangle\in V[[z]]$.
    \end{lemma}
\end{dang}
Proof: By $T-$covariance we have $[T,a_n]=-na_{n-1}\;\forall n\in\bb{Z}$. So we can prove by induction that $T^ka_n|0\rangle = (-1)^kn!/(n-k)!a_{n-k}|0\rangle$. Now, since $a(z)$ is a quantum field, we have $a_n|0\rangle=0$ whenever $n\geq N$ for a fixed $N\geq0$ depending on $|0\rangle$. In particular, this means that for any $k, 1\leq k\leq N$:
\begin{equation}
    0=T^ka_N|0\rangle = (-1)^k\frac{N!}{(N-k)!}a_{N-k}|0\rangle
\end{equation}
and $a_{N-k}|0\rangle=0\;(1\leq k\leq N)$. Therefore $a_n|0\rangle=0\;\forall n\geq0$ which is equivalent to $a(z)|0\rangle\in V[[z]]$.

\begin{dang}
    \begin{notation}
        From now on we will fix a vertex algebra $(V,|0\rangle, T,\mathcal{F})$. In this context we will write $\widetilde{\mathcal{F}}$ for the vector space of all $T$-covariant $\text{End}V-$valued quantum fields. By the lemma~(\ref{Lemma: a(z)|0> only have positive powers of z}), we are allowed to define a linear map $s:\Tilde{\mathcal{F}}\rightarrow V:a(z)\mapsto a(z)|0\rangle |_{z=0}$.
    \end{notation}
\end{dang}
\begin{dang}
    \begin{lemma}\label{Lemma: s(a(z))=0 implies a(z)|0>=0}\cite{de2006finite}
        Let $a(z)$ be a $T-$covariant quantum field and write $s(a(z))=a$. Then $a(z)|0\rangle=e^{zT}a$. In particular, this shows that if $a=0$, then $a(z)|0\rangle=0$.
    \end{lemma}
\end{dang}
Proof: By $T-$covariance we have $Ta_{(n)}|0\rangle=[T,a_{(n)}]|0\rangle=-na_{n-1}|0\rangle$ and by induction 
\begin{equation}
    T^ka_{(n)}|0\rangle =(-1)^kn(n-1)\cdots(n-k+1)a_{(n-k)}|0\rangle.
\end{equation}
In particular $T^ka_{(-1)}|0\rangle=k!a_{(-k-1)}|0\rangle\Rightarrow a_{(-k-1)}|0\rangle=T^ka_{(-1)}|0\rangle/k!=T^ka/k!$. By lemma~(\ref{Lemma: a(z)|0> only have positive powers of z}) we already know that $a_{(n)}|0\rangle=0\;\forall n\geq0$. Therefore $a(z)|0\rangle=\sum_{k\geq0}T^ka/k!=e^{zT}a$.\par
Let $\bb{C}[[z,w]]$ be the ring of Taylor series in two variables and $\bb{C}((z,w))$ its fraction field. We have the following natural homomorphism
\begin{equation}
    r_{z,w}:\bb{C}[[z,w]]\rightarrow\bb{C}((z))((w)):\sum_{m,n\geq0}a_{m,n}z^mw^n\mapsto \sum_{n\geq0}\left(\sum_{m\geq0}a_{m,n}z^m\right)w^n.
\end{equation}
Then, by the universal property of the fraction field, we have a homomorphism $i_{z,w}$ such that the following diagram commutes:
\begin{center}
    \begin{tikzcd}
       \bb{C}[[z,w]] \arrow[r, "r_{z,w}"] \arrow[d, hook]  &\bb{C}((z))((w)) \\
       \bb{C}((z,w))\arrow[ru, swap, "i_{z,w}"]
    \end{tikzcd}
\end{center}
Now, suppose that $D:\bb{C}[[z,w]]\rightarrow \bb{C}[[z,w]]$ is a derivation. It can be shown that there exists a unique extension of $D$ to a derivation of $\bb{C}((z,w))$ (which will denote by $D$ again). The formula is the usual one:
\begin{equation}
    D\left(\frac{f}{g}\right)= \frac{(Df)g-f(Dg)}{g^2}.
\end{equation}
Now, let $D=\partial_z$. Since $\partial_z\circ r_{z,w}=r_{z,w}\circ \partial_z$, we see that $\partial_z\circ i_{z,w}=i_{z,w}\circ\partial_z$ by the definition of the extension of $\partial_z$ to $ \bb{C}((z,w))$. In the same way we can see that $i_{z,w}$ commutes with multiplication by elements of $\bb{C}[[z,w]]$.\par
As an example, we have
\begin{equation}
    i_{z,w}\frac{1}{z-w}=\sum_{k\geq0}z^{-k-1}w^k
\end{equation}
because $(z-w)\sum_{k\geq0}z^{-k-1}w^k=1$ in $\bb{C}((z))((w))$. Applying $\partial_z^l$ to both sides of this equation we obtain
\begin{equation}
\begin{split}
     (-1)^ll!i_{z,w}\frac{1}{(z-w)^{l+1}}&=\sum_{k\geq0} (-1)^l\frac{(k+l)!}{k!}z^{-k-l-1}w^k\Rightarrow \\
     i_{z,w}\frac{1}{(z-w)^{l+1}}&=\sum_{k\geq0}\binom{k+l}{k}z^{-l-1-k}w^k.
     \end{split}
\end{equation}
Analytically, we can interpret $i_{z,w}(1/(z-w)^l)$ as the expansion of the series $1/(z-w)^l$ in the domain $|z|>|w|$. It's important to note that the diagram
\begin{center}
    \begin{tikzcd}
       \bb{C}((z,w)) \arrow[r, "i_{z,w}"] \arrow[d, "i_{w,z}"] &  \bb{C}((z))((w)) \arrow[d, hook]\\
       \bb{C}((w))((z)) \arrow[r, hook] & \bb{C}[[z^\pm,w^\pm]]
    \end{tikzcd}
\end{center}
is not commutative. In fact, we have the following important identity
\begin{equation}\label{Equation: formal delta distribution identity}
    i_{z,w}\frac{1}{z-w}-i_{w,z}\frac{1}{z-w}=\sum_{j\in\bb{Z}}z^{-j-1}w^j.
\end{equation}
There is a special notation for the series appearing on the right side of this equation.
\begin{dang}
    \begin{notation}
        We write $\delta(z,w)=\sum_{z,w}z^{-j-1}w^j$ and we call it the \textit{formal delta distribution}.
    \end{notation}
\end{dang}
\begin{dang}
    \begin{lemma}\label{Lemma: formal delta distribution identities}\cite{kac2013bombay}
        The formal delta distribution satisfies the following identities:
        \begin{enumerate}
            \item[(1)] $(z-w)^k\partial_w^j\delta(z,w)=0$ for $k>j$.
            \item[(2)] $(z-w)\partial_w^j\delta(z,w)=j\partial_w^{j-1}\delta(z,w)$.
        \end{enumerate}
    \end{lemma}
\end{dang}
Proof: Using equation~(\ref{Equation: formal delta distribution identity}) we see that
\begin{equation}
    \partial_w^j\delta(z,w)=j!\left(i_{z,w}\frac{1}{(z-w)^{j+1}}-i_{w,z}\frac{1}{(z-w)^{j+1}}\right).
\end{equation}
In particular, it's clear from this expression that $(z-w)^k\partial_w^j\delta(z,w)=0$ for $k>j$. Now
\begin{equation}
    \begin{split}
        (z-w) \partial_w^j\delta(z,w)&= j!\left(i_{z,w}\frac{z-w}{(z-w)^{j+1}}-i_{w,z}\frac{z-w}{(z-w)^{j+1}}\right)\\
        &=j \partial_w^{j-1}\delta(z,w)
    \end{split}
\end{equation}
proving the second identity.\par
As a consequence of the previous lemma, we have $z\delta(z,w)=w\delta(z,w)\Rightarrow f(z)\delta(z,w)=f(w)\delta(z,w)$ for any $U-$valued formal distribution $f$.\par
Now we can introduce the notion of $n$th product of $\text{End}V-$valued quantum fields. Let $U$ be a vector superspace and define the linear map $\text{Res}_z:U[[z^\pm]]\rightarrow U: a(z)\mapsto a_0$, i.e., extraction of the coefficient of $z^{-1}$. Given two quantum fields $a(z), b(z)$, their $n$th product is defined to be
\begin{equation}\label{Equation: nth product of quantum fields}
    a(w)_{(n)}b(w)=\text{Res}_z (i_{z,w}(z-w)^na(z)b(w)-p(a,b)i_{w,z}(z-w)^nb(w)a(z)).
\end{equation}
\begin{dang}
  \begin{lemma}\label{Lemma: the nth product of quamtum fields is a quantum field}\cite{de2006finite}
      Let $a(z), b(z)$ be two quantum fields on $V$. Then, their $n$th product $a(z)_{(n)}b(z)$ is again a quantum field $\forall n\in \bb{Z}$. 
  \end{lemma}
\end{dang}
Proof: Let $v\in V$, since $a(z)$ and $b(w)$ are quantum fields, we have $b(w)v\in V((w))$ and $a(z)b(w)v\in V((z))((w))$. But $i_{z,w}(z-w)^n\in \bb{C}((z))[[w]]$, so the product $i_{z,w}(z-w)^na(z)b(w)v$ is well defined as an element of $V((z))((w))$, which implies that
\[
\text{Res}_zi_{z,w}(z-w)^na(z)b(w)v \in V((w)).
\]
Here we have used the inclusion
\[V((z))((w))\hookrightarrow V[[z^\pm,w^\pm]]\cong V[[w^\pm]][[z^\pm]]\]
 to apply $\text{Res}_z$. Similarly with the second term, we get $\text{Res}_zi_{w,z}(z-w)^nb(w)a(z)\in V((w))$. Therefore $a(w)_{(n)}b(w)v\in V((w))\;\forall v\in V$ and $a(w)_{(n)}b(w)$ must be a quantum field. 
\begin{dang}
    \begin{lemma}\label{Lemma: nth products of T-covariants is T-covariant}\cite{de2006finite}
        Given two $T-$covariant quantum fields $a(z), b(z)$, their $n$th product $a(z)_{(n)}b(z)$ is also a $T-$covariant quantum field.
    \end{lemma}
\end{dang}
Proof: By lemma~(\ref{Lemma: the nth product of quamtum fields is a quantum field}), we just need to verify that $a(z)_{(n)}b(z)$ is $T-$covariant. Now, using the fact that $a(z), b(z)$ are $T-$covariant quantum fields, we get
\begin{equation}
    Ta(z)b(w) = a(z)b(w)T+\partial_za(z)b(w) + a(z)\partial_wb(w).
\end{equation}
Therefore
\begin{equation}
\begin{split}
    &[T,a(w)_{(n)}b(w)] \\
    &\quad= \text{Res}_z( i_{z,w}(z-w)^n\partial_za(z)b(w) + i_{z,w}(z-w)^na(z)\partial_wb(w)\\
    &\quad\quad-p(a,b)i_{w,z}(z-w)^n\partial_wb(w)a(z)-p(a,b)i_{w,z}(z-w)^nb(w)\partial_za(z)).
\end{split}
\end{equation}
But 
\begin{equation}
\begin{split}
     &\partial_zi_{z,w}(z-w)^na(z)b(w) + \partial_wi_{w,z}(z-w)^na(z)b(w)\\
     &= i_{z,w}(z-w)^n\partial_za(z)b(w) + i_{z,w}(z-w)^na(z)\partial_wb(w)
\end{split}
\end{equation}
and
\begin{equation}
\begin{split}
    &\partial_zi_{w,z}(z-w)^nb(w)a(z) + \partial_wi_{w,z}(z-w)^nb(w)a(z)\\
    &=i_{w,z}(z-w)^nb(w)\partial_za(z) + i_{w,z}(z-w)^n\partial_wb(w)a(z).
\end{split} 
\end{equation}
Using the fact that $\text{Res}_z\circ\partial_w=\partial_w\circ\text{Res}_z, 
\;\text{Res}_z\circ \partial_z =0$, and the two previous equations, we get $[T,a(w)_{(n)}b(w)]=\partial_w(a(w)_{(n)}b(w))$.\par
Another important property of $n$th products is the following.
\begin{dang}
    \begin{lemma}\label{Lemma: d(a(w)nb(w))=da(w)nb(w)+a(w)ndb(w)}
        $\partial_w(a(w)_{(n)}b(w))=(\partial_wa(w))_{(n)}b(w)+a(w)_{(n)}\partial_w(b(w))$.
    \end{lemma}
\end{dang}
Proof: First we have
\begin{equation}
    \begin{split}
        \partial_w(a(w)_{(n)}b(w)) = \text{Res}_z( -ni_{z,w}(z-w)^{n-1}a(z)b(w) +i_{z,w}(z-w)^na(z)\partial_wb(w)\\+p(a,b)ni_{w,z}(z-w)^{n-1}b(w)a(z)-p(a,b)i_{w,z}(z-w)^{n}\partial_wb(w)a(z)).
    \end{split}
\end{equation}
But 
\begin{equation}
\begin{split}
    -ni_{z,w}(z-w)^{n-1}a(z)b(w) &= -\partial_zi_{z,w}(z-w)^na(z)b(w)+i_{z,w}(z-w)^n\partial_za(z)b(w)\quad\text{ and}\\
    ni_{w,z}(z-w)^{n-1}b(w)a(z)&=\partial_zi_{w,z}(z-w)^{n}b(w)a(z)-i_{w,z}(z-w)^{n}b(w)\partial_za(z).
\end{split}
\end{equation}
So using that $\text{Res}_z\circ \partial_z=0$ we get the result.
\par
Now we enunciate ``Dong's lemma''.
\begin{dang}
    \begin{lemma}\label{Lemma: Dong's lemma}\cite{de2006finite}
        Given a triple $(a(z),b(z),c(z))$ of pairwise local quantum fields , then $(a(z)_{(n)}b(z), c(z))$ is a local pair of quantum field, for any $n\in\bb{Z}$.
    \end{lemma}
\end{dang}
Proof: Fix $n\in\bb{Z}$. By hypothesis, there is an integer $r\geq |n|$ such that
\begin{equation}
   \begin{split}
       &(z-w)^r[a(z),b(w)]=0,\\
        &(z-x)^r[a(z),c(x)]=0,\\
         &(w-x)^r[b(w),c(x)]=0.
   \end{split}
\end{equation}
Now let $N>4r$. By the binomial formula we have
\begin{equation}
    (x-w)^N=(x-w)^r\sum_{s\geq0}\binom{N-r}{s}(x-z)^{N-r-s}(z-w)^s.
\end{equation}
So
\begin{equation}\label{Equaiton: Dong's lemma equation}
\begin{split}
     &(x-w)^N[c(x),a(w)_{(n)}b(w)]\\
     &=\text{Res}_z(x-w)^{r}\sum_{s\geq0}\binom{N-r}{s}(x-z)^{N-r-s}(z-w)^s\times\\
     &\quad\times(i_{z,w}(z-w)^n[c(x),a(z)b(w)]-p(a,b)i_{w,z}(z-w)^n[c(x),b(w)a(z)]).
\end{split}
\end{equation}
But 
\begin{equation}
    [c(x),a(z)b(w)]-p(a,b)[c(x),b(w)a(z)]=c(x)[a(z),b(w)]+(-1)^{p(c)}[b(w),a(z)]c(x)
\end{equation}
and consequently 
\begin{equation}
\begin{split}
    &(z-w)^s(i_{z,w}(z-w)^n[c(x),a(z)b(w)]-p(a,b)i_{w,z}(z-w)^n[c(x),b(w)a(z)])\\
    &\quad=(z-w)^{n+r}(z-w)^{s-r}(c(x)[a(z),b(w)]+(-1)^{p(c)}[b(w),a(z)]c(x))\\
    &\quad=0
\end{split}
\end{equation}
whenever $s>2r$. Moreover, we can write
\begin{equation}
\begin{split}
     &[c(x),a(z)b(w)]=[c(x),a(z)]b(w)+p(a,c)a(z)[c(x),b(w)]\quad\text{ and}\\
     &\quad [c(x),b(w)a(z)]=[c(x),b(w)]a(z)+p(b,c)b(w)[c(x),a(z)].
\end{split}
\end{equation}
Therefore, if $s\leq 2r\Rightarrow N-r-s>r$ and 
\begin{equation}
    (x-w)^r(x-z)^{N-r-s}[c(x),a(z)b(w)]=0=(x-w)^r(x-z)^{N-r-s}[c(x),b(w)a(z)].
\end{equation}
So, we conclude that all the terms in the equation~(\ref{Equaiton: Dong's lemma equation}) are zero, proving the locality of the pair $(c(x),a(z)_{(n)}b(z))$.
\begin{dang}
    \begin{lemma}\label{Lemma: s(a(z)_nb(z))=a_nb}\cite{de2006finite}
        Let $a(z), b(z)$ be quantum fields and $s(a(z))=a,\; s(b(z))=b$. Then $s(a(z)_{(n)}b(z))=a_{(n)}b\;\forall n\in\bb{Z}$.
    \end{lemma}
\end{dang}
Proof: By the equation~(\ref{Equation: nth product of quantum fields}), we have
\begin{equation}
\begin{split}
     &a(w)_{(n)}b(w)|0\rangle=\\
     &\quad\text{Res}_z  (i_{z,w}(z-w)^na(z)b(w)|0\rangle-p(a,b)i_{w,z}(z-w)^nb(w)a(z)|0\rangle).
\end{split}
\end{equation}
But $a(z)|0\rangle\in V[[z]]$ and $i_{w,z}(z-w)^n\in V((w))[[z]]\Rightarrow \text{Res}_zi_{w,z}(z-w)^nb(w)a(z)|0\rangle=0$. Therefore
\begin{equation}
    \begin{split}
        a(w)_{(n)}b(w)|0\rangle\big|_{w=0}&=\text{Res}_zi_{z,w}(z-w)^na(z)b(w)|0\rangle\big|_{w=0}\\
        &=\text{Res}_zz^na(z)b\\
        &=a_{(n)}b.
    \end{split}
\end{equation}
The next task will be to prove the \textit{state-field correspondence}. Before we proceed, let's make some observations. The first one is that $\partial_za(z)$ is $T-$covariant whenever $a(z)$ is $T-$covariant. Moreover, for any local pair $(a(z),b(z))$, $(\partial_za(z),b(z))$ is also a local pair. Now, let $\mathcal{F}^\prime$ be the smallest subspace of $\Tilde{\mathcal{F}}$ that contains $I_V$, is local with each element of $\mathcal{F}$, is closed under $\partial_z$ and $n$th products. We can construct $\mathcal{F}^\prime$ inductively as follows. Let $\mathcal{F}_0=\text{Span}\{\mathcal{F},I_V\}$ and for $k\geq0$
\begin{equation}
\mathcal{F}_{k+1}=\mathcal{F}_k+\sum_{n\geq0}\partial_z^n\mathcal{F}_k+\sum_{n\geq0}(\mathcal{F}_k)_{(n)}\mathcal{F}_k.
\end{equation}
The fact that the elements of $(\mathcal{F}_k)_{(n)}\mathcal{F}_k\;(k\geq0)$ are local with $\mathcal{F}$ is a consequence of the Dong's lemma~(\ref{Lemma: Dong's lemma}). So we put $\mathcal{F}^\prime=\bigcup_{k\geq0}\mathcal{F}_k$. Now, let $\Bar{\mathcal{F}}$ be the subspace of $\Tilde{\mathcal{F}}$ consisting of $\text{End}V-$quantum fields that are local with $\mathcal{F}$. By definition, we have the sequence of inclusions
\begin{equation}
    \mathcal{F}\subset\mathcal{F}^\prime\subset\Bar{\mathcal{F}}\subset\Tilde{\mathcal{F}}.
\end{equation}
\begin{dang}
    \begin{theorem}\label{Theorem: state field correspondence}\cite{de2006finite}
        Let $(V,T,|0\rangle,\mathcal{F})$ be a vertex algebra. Then the linear map $s:\mathcal{F}^\prime\rightarrow V$ is bijective.
    \end{theorem}
\end{dang}
Proof: Suppose that $a(z)$ is a $T-$covariant quantum field that is local with each element of $\mathcal{F}$. Then by Dong's lemma~(\ref{Lemma: Dong's lemma}), $a(z)$ is local with $n$th products of elements of $\mathcal{F}$ and its derivatives. This proves that the pair $(a(z),b(z))$, for $b(z)\in\mathcal{F}^\prime,a(z)\in\Bar{\mathcal{F}}$, is local. Now, by the completeness axiom, the map $s:\mathcal{F}^\prime\rightarrow V$ is surjective. This follows by the identity
\begin{equation}
    \begin{split}
        s(a^{j_1}(z)_{(n_1)}(a^{j_2}(z)_{(n_2)}\cdots (a^{j_k}(z)_{(n_k)}I_V)\cdots) &= a_{n_1}^{j_1}\cdots a_{n_k}^{j_k}s(I_V)\quad(\text{by lemma~(\ref{Lemma: s(a(z)_nb(z))=a_nb})})\\
        &=a_{n_1}^{j_1}\cdots a_{n_k}^{j_k}|0\rangle.
    \end{split}
\end{equation}
In addition, we claim that the map $s:\Bar{\mathcal{F}}\rightarrow V$ is injective. In fact, let and $a(z)\in\Bar{\mathcal{F}}$ and suppose that $s(a(z))=0$. Now, suppose that $b(z)\in\mathcal{F}^\prime$. Previously, we have shown that the pair $(a(z),b(z))$ is local. So there is an $N\geq0$ such that
\begin{equation}
\begin{split}
     &(z-w)^Na(z)b(w)=p(a,b)(z-w)^Nb(w)a(z)\Rightarrow\\
     &\;(z-w)^Na(z)b(w)|0\rangle=p(a,b)(z-w)^Nb(w)a(z)|0\rangle\Rightarrow\\
     &\;\;(z-w)^Na(z)b(w)|0\rangle=0\quad(\text{because of lemma~(\ref{Lemma: s(a(z))=0 implies a(z)|0>=0})})\Rightarrow\\
     &\quad (z-w)^Na(z)b(w)|0\rangle\big|_{w=0}= z^Na(z)b=0.
\end{split}
\end{equation}
Since $a(z)b\in V((z))$, we can cancel out $z^N$ in the last equation, to get $a(z)b=0\Rightarrow a_{(n)}b=0\;\forall n\in\bb{Z}$. Since $s:\mathcal{F}^\prime\rightarrow V$ is surjective, this shows that $a(z)=0$. Now we have the commutative diagram
\begin{center}
    \begin{tikzcd}
        \mathcal{F}^\prime \arrow[rr, hook]\arrow[dr, two heads, swap, "s|_{\mathcal{F}^\prime}"] &  & \Bar{\mathcal{F}}\arrow[ld, hook, "s|_{\Bar{\mathcal{F}}}"] \\
         & V &
    \end{tikzcd}
\end{center}
which implies that $ \mathcal{F}^\prime=\Bar{\mathcal{F}}$ and that $s:\mathcal{F}^\prime\rightarrow V$ is bijective.\par
Theorem~(\ref{Theorem: state field correspondence}) implies the existence of a linear map that associates to each $a\in V$ a quantum field $Y(a,z)\in\mathcal{F}^\prime$ such that $Y(a,z)|0\rangle\big|_{z=0}=a$. Although the definition of $Y(a,z)$ is constructive, we can develop techniques to do calculations with these quantum fields. For example, we have the \textit{Borcherds identity}:
\begin{equation}\label{equation: Borcherds identity}
\begin{split}
     i_{z,w}(z-w)^nY(a,z)Y(b,w)-p(a,b)i_{w,z}(z-w)^nY(b,w)Y(a,z)\\=\sum_{j\geq0}Y(a_{(n+j)}b,w)\partial_z^j\delta(z,w)/j!
\end{split}
\end{equation}
where $Y(a_{(n+j)}b,w)=Y(a,w)_{(n+j)}Y(b,w)$. To show this identity we need a lemma.
\begin{dang}
    \begin{lemma}\label{Lemma: operator product expansion}\cite{kac2013bombay}
        Let $D=D(z,w)$ be a local distribution, i.e. $(z-w)^ND=0$ for some $N\geq0$. Then
        \begin{equation}
            D = \sum_{j=0}^{N-1}c_j(w)\partial_z^j\delta(z,w)/j!
        \end{equation}
        where $c_j(w)=\text{Res}_z(z-w)^jD$. This equation is called \textit{operator product expansion} of the formal distribution $D$.
    \end{lemma}
\end{dang}
Proof: The proof follows by induction on $N$. Let $D=\sum_{n,m\in\bb{Z}}D_{n,m}z^nw^m$. Suppose that $(z-w)D=0\Rightarrow D_{n,m}=D_{n+1,m-1}$ and
\begin{equation}
    \begin{split}
        D &= \sum_{k\in\bb{Z}}\sum_{n+m=k}D_{n,m}z^nw^m\\
          &=\sum_{k\in\bb{Z}}D_{k,0}\sum_{n\in\bb{Z}}z^nw^{k-n}\\
          &=\sum_{k\in\bb{Z}}D_{k,0}w^{k+1}\delta(z,w)\\
          &=c_0(w)\delta(z,w).
    \end{split}
\end{equation}
Also, we have $\text{Res}_z D=c_0(w)\text{Res}_z\delta(z,w)=c_0(w)$. This proves the base case $N=1$. Now, suppose that the statement of the lemma is true for some $N>0$ and suppose that $(z-w)^{N+1}D=0$. Define $D^\prime = (z-w)^ND\Rightarrow (z-w)D^\prime=0$ and $D^\prime = c_N(w)\delta(z,w)$ by the first part of the proof. Now, define $D^{\prime\prime}= D - c_N(w)\partial_w^N\delta(z,w)/N!$. Then, by lemma~(\ref{Lemma: formal delta distribution identities}) we get
\begin{equation}
\begin{split}
    (z-w)^ND^{\prime\prime} &= (z-w)^ND-(z-w)^N c_N(w)\partial_w^N\delta(z,w)/N!\\
    &=c_N(w)\delta(z,w)-c_N(w)\delta(z,w)\\
    &=0.
\end{split}
\end{equation}
By the induction hypothesis, we have $D^{\prime\prime}=\sum_{j=0}^{N-1}c_j(w)\partial_w^j\delta(z,w)/j!$ where $c_j(w)=\text{Res}_z(z-w)^jD^{\prime\prime} = \text{Res}_z(z-w)^jD$ because of the lemma~(\ref{Lemma: formal delta distribution identities}) and the fact $\text{Res}_z\partial_w^k\delta(z,w)=0\;\forall k>0$. Therefore, we get $D=\sum_{j=0}^Nc_j(w)\partial_w^j\delta(z,w)/j!$ where $c_j(w)=\text{Res}_z(z-w)^jD$.\par
Now, we are ready to prove the Borcherds identity. First, the pair $(Y(a,z),Y(b,z))$ is local, so there is a nonnegative integer $N\geq0$ such that $(z-w)^N[Y(a,z),Y(b,w)]=0$. So, multiplying the left hand side of the equation~(\ref{equation: Borcherds identity}) by $(z-w)^{N+|n|}$ we get zero. Therefore, the expression 
\begin{equation}
    i_{z,w}(z-w)^nY(a,z)Y(b,w)-p(a,b)i_{w,z}(z-w)^nY(b,w)Y(a,z)
\end{equation}
is a local distribution. Moreover, for $j\geq0$
\begin{equation}
\begin{split}
    \text{Res}_z(z-w)^j(i_{z,w}(z-w)^nY(a,z)Y(b,w)-p(a,b)i_{w,z}(z-w)^nY(b,w)Y(a,z))\\
    =Y(a,w)_{(n+j)}Y(b,w)
\end{split}
\end{equation}
and by applying the lemma~(\ref{Lemma: operator product expansion}) we get equation~(\ref{equation: Borcherds identity}).\par
\begin{dang}
    \begin{definition}\label{Definition: normal ordered product between fields}\cite{de2006finite}
        Given two $\text{End}V-$valued quantum fields $a(z), b(z)$ we define the \textit{normal ordered product} $:a(z)b(z):=a(z)_+b(z)+a(z)b(z)_-$ where 
        \begin{equation}
            a(z)_+=\sum_{n<0}a_{(n)}z^{-n-1}\;\;\text{and}\;\;a(z)_-=\sum_{n\geq0}a_{(n)}z^{-n-1}.
        \end{equation}
    \end{definition}
\end{dang}
Concerning the operators $(\cdot)_\pm:\text{End}V[[z^\pm]]\rightarrow\text{End}V[[z^\pm]$, we have the useful identities:
\begin{equation}
    \text{Res}_za(z)i_{z,w}\frac{1}{z-w}=a(w)_+\;\;\text{and}\;\;\text{Res}_za(z)i_{w,z}\frac{1}{z-w}=-a(w)_-.
\end{equation}
Both identities can be verified by direct calculation. For example
\begin{equation}
    \begin{split}
         \text{Res}_za(z)i_{z,w}\frac{1}{z-w}&=\text{Res}_z\sum_{n\in\bb{Z},k\geq0}a_{(n)}z^{-k-1}w^k\\
         &=\sum_{k\geq0,n\in\bb{Z}}\sum_{n+k=-1}a_{(n)}w^k=\sum_{n<0}a_{(n)}w^{-n-1}.      
    \end{split}
\end{equation}
Now, using the definition of $n$th products of fields~(\ref{Equation: nth product of quantum fields}) for $n=-1$ we get the important identity $:Y(a,z)Y(b,z):=Y(a,z)_{(-1)}Y(b,z)=Y(a_{(-1)}b,z)$. So we define the normal ordered product in $V$ as the bilinear map $V\otimes V\rightarrow V:(a,b)\mapsto\; :ab:= a_{(-1)}b$. Sometimes it can be useful to explicitly open the definition of the normal ordered product between quantum fields. In general we have the identity
\begin{equation}\label{Equation: normal ordered product between quantum fields}
    :a(z)b(z):=\sum_{n\in\bb{Z}}\sum_{m<0}(a_{(m)}b_{(n-m-1)}+b_{(n+m)}a_{(-m-1)}) z^{-n-1}.
\end{equation}

\section{Calculating in vertex algebra}
Let $a(z)\in U((z))$ be a $U-$valued formal distribution. We define a map $U((z))\rightarrow U[[\lambda]]$ by the formula $F_z^\lambda a(z)=\text{Res}_ze^{\lambda z}a(z)$, which is called the \textit{formal Fourier transform}. We can verify by direct computation that
\begin{equation}
      F_z^\lambda a(z) = \text{Res}_z\sum_{j\geq0}\sum_{n\in\bb{Z}}\frac{\lambda^j}{j!}a_{(n)}z^{j-n-1} =\sum_{j\geq0}\frac{\lambda^j}{j!}a_{(j)}.
\end{equation}
\begin{dang}
    \begin{lemma}\label{Lemma: formal fourier transform identities}\cite{de2006finite}
        Let $a(z)\in U((z))$. The formal Fourier transform satisfies the following identities:
        \begin{enumerate}
            \item[(1)] $F_z^\lambda(e^{zT}a(z))=F_z^{\lambda+T}a(z)$.
            \item[(2)] $F_z^\lambda a(-z)=-F_z^{-\lambda}a(z)$.
            \item[(3)] $F_z^\lambda\partial_w^n\delta(z,w)=e^{\lambda w}\lambda^n$.
        \end{enumerate}
    \end{lemma}
\end{dang}
Proof: The first identity is immediate by the definition. For the second identity
\begin{equation}
    \begin{split}
         F_z^\lambda a(z) &= \text{Res}_z\sum_{j\geq0}\sum_{n\in\bb{Z}}\frac{\lambda^j}{j!}a_{(n)}(-1)^{-n-1}z^{j-n-1}\\
         &=\sum_{j\geq0}\frac{\lambda^j}{j!}a_{(j)}(-1)^{-j-1}\\
         &=-\sum_{j\geq0}\frac{(-\lambda)^j}{j!}a_{(j)}\\
         &=-F_z^{-\lambda}a(z).
    \end{split}
\end{equation}
Now, for the third identity
\begin{equation}
    \begin{split}
        F_z^\lambda\partial_w^n\delta(z,w)&=\text{Res}_ze^{\lambda z}\partial_w^n\delta(z,w)\\
        &=(-1)^n\text{Res}_ze^{\lambda z}\partial_z^n\delta(z,w)\\
        &=\text{Res}_z(\partial_w^ne^{\lambda z})\delta(z,w)\\
        &=\text{Res}_z(\partial_w^ne^{\lambda w})\delta(z,w)\\
        &=\partial_w^ne^{\lambda w}\text{Res}_z\delta(z,w)\\
        &=\partial_w^ne^{\lambda w}\cdot 1= e^{\lambda w}\lambda^n.
    \end{split}
\end{equation}
Now, let $(V,T,|0\rangle,\mathcal{F})$ be a vertex algebra. We define a bilinear map $[\cdot_\lambda\cdot]:V\otimes V\rightarrow\bb{C}[[\lambda]]\otimes V$ called \textit{lambda-bracket} as follows. For $a,b\in V$ we put $[a_\lambda b]=F_z^\lambda Y(a,z)b$.
\begin{dang}
    \begin{theorem}\label{Theorem: lambda bracket identities}\cite{de2006finite}
    Let $(V,T,|0\rangle,\mathcal{F})$ be a vertex algebra. The lambda bracket satisfies the following identities:
    \begin{enumerate}
        \item[(1)] (sesquilinearity) $[Ta_\lambda b]=-\lambda[a_\lambda b]$ and $T[a_\lambda b]=[Ta_\lambda b]+[a_\lambda Tb]$.
        \item[(2)] (skewsymmetry) $[b_\lambda a]=-p(a,b)[a_{-\lambda-T}b]$.
        \item[(3)] (noncommutative Wick formula) 
        \begin{equation}
            [a_{\lambda}:bc:]=:[a_\lambda b]c:+p(a,b):b[a_\lambda c]:+\int_0^\lambda[[a_\lambda b]_\mu c]d\mu.
        \end{equation}
    \end{enumerate}
    \end{theorem}
\end{dang}
Before we prove this theorem, we need a few lemmas.
\begin{dang}
    \begin{lemma}\label{Lemma: T(anb)=Tanb+anTb and Tan=-na(n-1)}\cite{de2006finite}
        Let $(V,T,|0\rangle,\mathcal{F})$ be a vertex algebra. Then we have the identities
        $(Ta)_{(n)}=-na_{n-1}$ and $T(a_{(n)}b)=(Ta)_{(n)}b+a_{(n)}(Tb)$.
    \end{lemma}
\end{dang}
Proof: Since $\partial_zY(a,z)|0\rangle\big|_{z=0}=TY(a,z)|0\rangle\big|_{z=0}=Ta$, we have $Y(Ta,z)=\partial_zY(a,z)$ because of the state-field correspondence. This implies the identity $(Ta)_{(n)}=-na_{n-1}$. Now, by lemma~(\ref{Lemma: s(a(z))=0 implies a(z)|0>=0}) we have $e^{zT}(a_{(n)}b)=Y(a,z)_{(n)}Y(b,z)|0\rangle$ and
\begin{equation}
\begin{split}
      Te^{zT}(a_{(n)}b)=(\partial_zY(a,z))_{(n)}Y(b,z)|0\rangle+Y(a,z)_{(n)}\partial_zY(b,z)|0\rangle\\
      =Y(Ta,z)_{(n)}Y(b,z)|0\rangle+Y(a,z)_{(n)}Y(Tb,z)|0\rangle
\end{split}
\end{equation}
by lemma~(\ref{Lemma: d(a(w)nb(w))=da(w)nb(w)+a(w)ndb(w)}). By putting $z=0$ in the previous equation we get the result.
\begin{dang}
    \begin{corollary}\label{Corollary: T acts as a derivaiton of the normal ordered product}\cite{de2006finite}
        Let $(V,|0\rangle,T,\mathcal{F})$ be a vertex algebra. Then $T(:ab:)=:(Ta)b:+:a(Tb):$.
    \end{corollary}
\end{dang}
\begin{dang}
    \begin{lemma}\label{Lemma: cojugation by the action of infinitesimal translations}
        Let $X\in U[[z^\pm]]$ and $T\in\text{End}U$. Then we have the identity $e^{wT}Xe^{-Tw}=e^{\text{ad}(wT)}X$.
    \end{lemma}
\end{dang}
Proof: By induction, it can be verified that 
\begin{equation}
    \text{ad}^k(T)=\sum_{j=0}^k(-1)^j\binom{k}{j}T^{k-j}XT^j.
\end{equation}
Therefore
\begin{equation}
    \begin{split}
        e^{wT}Xe^{-Tw}&=\sum_{k\geq0}\sum_{j=0}^k\frac{w^{k-j}}{(k-j)!}T^{k-j}XT^j\frac{(-w)^j}{j!}\\
        &=\sum_{k\geq0}\sum_{j=0}^k\frac{(-1)^j}{(k-j)!j!}T^{k-j}XT^jw^k\\
        &=\sum_{k\geq0}\frac{1}{k!}\sum_{j=0}^k(-1)^j\binom{k}{j}T^{k-j}XT^jw^k\\
        &=\sum_{k\geq0}\frac{w^k}{k!}\text{ad}^k(T)X.
    \end{split}
\end{equation}
\begin{dang}
    \begin{lemma}\label{Lemma: formal taylor lemma}
        Let $f(z)\in U[[z^\pm]]$ be a formal distribution. Then $e^{w\partial_z}f(z)=f(z+w)$, where the terms $(z+w)^{-1}$ are expanded in $\bb{C}((z))((w))$, i.e., with positive powers of $w$.
    \end{lemma}
\end{dang}
Proof: Let $f(z)=\sum_{n\in\bb{Z}}f_nz^n$. Then 
\begin{equation}
    \partial_z^kf(z)=\sum_{n\in\bb{Z}}f_nn(n-1)\cdots(n-k+1)z^{n-k}
\end{equation}
and
\begin{equation}
\begin{split}
    e^{w\partial_z}f(z)&=\sum_{k\geq0}\frac{w^k}{k!}\partial_z^kf(z)\\
    &=\sum_{n\in\bb{Z}}f_n\sum_{k\geq0}\frac{n(n-1)\cdots(n-k+1)}{k!}z^{n-k}w^k\\
    &=\sum_{n\in\bb{Z}}f_ni_{z,w}(z+w)^n\\
    &=f(z+w).
\end{split}
\end{equation}
\begin{dang}
    \begin{lemma}\label{Lemma: skewsymmetry for fields}\cite{de2006finite}
        Let $(V,T,|0\rangle,\mathcal{F})$ be a vertex algebra. Then we have the \textit{skewsymmetry} relation $Y(a,z)b=p(a,b)e^{zT}Y(b,-z)a\;\;\forall a,b\in V$.
    \end{lemma}
\end{dang}
Proof: First, for $a\in V$ we have the sequence of identities
\begin{equation}\label{Equation: conjugation of quantum fields by the action of infinitesimal translations}
    \begin{split}
        e^{wT}Y(a,z)e^{-Tw}&= e^{\text{ad}(wT)}Y(a,z)\quad(\text{by lemma~(\ref{Lemma: cojugation by the action of infinitesimal translations})})\\
        &=e^{w\partial_z}Y(a,z)\quad(\text{by translation covariance})\\
        &=Y(a,z+w)\quad(\text{by lemma~(\ref{Lemma: formal taylor lemma})}).
    \end{split}
\end{equation}
Now let $b\in V$, by locality there is a integer $N\geq0$ such that $(z-w)^NY(a,z)Y(b,w)=p(a,b)(z-w)^NY(b,w)Y(a,z)$. Applying  both sides we to the vacuum, we get
\begin{equation}\label{equation: skewsymmetry relation}
\begin{split}
      (z-w)^NY(a,z)e^{wT}b&=p(a,b)(z-w)^NY(b,w)e^{zT}a\\
      &=p(a,b)(z-w)^Ne^{-zT}Y(b,w-z)a\quad(\text{by equation~(\ref{Equation: conjugation of quantum fields by the action of infinitesimal translations})}),
\end{split}
\end{equation}
because of Lemma~(\ref{Lemma: s(a(z))=0 implies a(z)|0>=0}). We remark that the terms $(z+w)^{-1}$ appearing in $Y(a,w-z)$ are expanded in $\bb{C}((z))((w))$. Now, with $N$ big enough, there is no negative power of $(z-w)$ in the equation~(\ref{equation: skewsymmetry relation}), so we get an identity in $V[[z]][[w]]$. Therefore, we can put $w=0$ in the equation~(\ref{equation: skewsymmetry relation}) and cancel $z^N$ on both sides to get the skewsymmetry relation.\par
Now we can give the proof of Theorem~(\ref{Theorem: lambda bracket identities}). Lemma~(\ref{Lemma: T(anb)=Tanb+anTb and Tan=-na(n-1)}) together with the relation $[a_\lambda b]=\sum_{j\geq0}\lambda^ja_{(j)}b/j!$ immediately implies the sesquilinearity identities of the theorem~(\ref{Theorem: lambda bracket identities}). Now, for the skewsymmetry we use Lemma~(\ref{Lemma: skewsymmetry for fields}) to obtain
\begin{equation}
    \begin{split}
        F_z^\lambda Y(a,z)b &= p(a,b)F_z^\lambda e^{zT}Y(b,-z)a\\
        &=-p(a,b)F_z^{-\lambda}e^{-zT}Y(b,z)a\quad(\text{by lemma~(\ref{Lemma: formal fourier transform identities})})\\
        &=-p(a,b)F_z^{-\lambda-T}Y(b,z)a\\
        &=-p(a,b)[b_{-\lambda-T}a].
    \end{split}
\end{equation}
It remains to prove the noncommutative Wick formula. The Borcherds identity~(\ref{equation: Borcherds identity}) for $n=0$ yields the commutator formula
\begin{equation}
\begin{split}
     &[Y(a,z),Y(b,w)]=\sum_{j\geq0}Y(a_{(j)}b,w)\partial_w^j\delta(z,w)/j!\Rightarrow\\
     &Y(a,z)Y(b,w)c-p(a,b)Y(b,w)Y(a,z)c=\sum_{j\geq0}Y(a_{(j)}b,w)c\partial_w^j\delta(z,w)/j!\;\;\forall c\in V.
\end{split}
\end{equation}
Applying $F_z^\lambda$ to both sides of this equation and using the third identity of lemma~(\ref{Lemma: formal fourier transform identities}), we get
\begin{equation}
\begin{split}
    [a_\lambda Y(b,w)c]-p(a,b)Y(b,w)[a_\lambda c]&=\sum_{j\geq0}Y(a_{(j)}b,w)c\lambda^je^{w\lambda}/j!\\
    &=e^{\lambda w}Y\left(\sum_{j\geq0}\frac{\lambda^j}{j!}a_{(j)}b\;,w\right)c\\
    &=e^{w\lambda}Y([a_\lambda b],w)c.
\end{split}
\end{equation}
Now, by applying $\text{Res}_z w^{-1}$ on both sides of this equation we obtain
\begin{equation}
\begin{split}
    [a_\lambda :bc:]-p(a,b):b[a_\lambda c]: &= \text{Res}_ww^{-1}e^{\lambda w}Y([a_\lambda b],w)c\\
    &=\text{Res}_w\left(\int_0^\lambda e^{\mu w}d\mu + w^{-1}\right)Y([a_\lambda b],w)c\\
    &=\int_0^\lambda \text{Res}_we^{\mu w}Y([a_\lambda b],w)c\;d\mu+\text{Res}_ww^{-1}Y([a_\lambda b],w)c\\
    &=\int_0^\lambda [[a_\lambda b]_\mu c]d\mu +:[a_\lambda b]c:
\end{split}
\end{equation}
which finishes the proof of the noncommutative Wick formula.

\section{Fermionic vertex algebras}
\begin{dang}
\begin{definition}\cite{kac2013bombay}
    Let $V=V_{\Bar{0}}\oplus V_{\Bar{1}}$ be a vector superspace. A \textit{skew-supersymmetric} bilinear form on $V$ is a bilinear map $\langle\cdot,\cdot\rangle: V\otimes V\rightarrow\bb{C}$ satisfying $\langle a,b\rangle=-p(a,b)\langle b,a\rangle$ for homogeneous elements $a,b\in V$.
\end{definition}
\end{dang}
In this section we will give our first example of a vertex algebras. Let $A=A_{\Bar{0}}\oplus A_{\Bar{1}}$ be a vector superspace with a skew-supersymmetric bilinear form $\langle\cdot,\cdot\rangle$, and write $A[t^{\pm}]=\bb{C}[t^{\pm}]\otimes A$. Now we define the \textit{Clifford affinization} of $A$ to be the Lie superalgebra $ \hat{A}=A[t^{\pm}]\oplus\bb{C}K$ with commutator given by
\begin{equation}\label{Equation: commutator in the clifford afinization}
   [at^m,bt^n]=\langle a,b\rangle\delta_{m,-n-1}K\quad\text{and}\quad [K,\cdot\;]=0.
\end{equation}
The parity on $\hat{A}$ is defined by putting $p(t)=0$ (i.e $\bb{C}[t^\pm]$ is defined to be a purely even superalgebra) and $p(K)=0$. So if $a$ homogeneous in $A$, we have $p(at^m)=p(a)\;\forall m\in\bb{Z}$. Define a derivation $D=-d/dt$ of $\hat{A}$ by putting $D(at^n)=-nat^{n-1}$ and $D(K)=0$. We claim that this in fact defines a derivation of $\hat{A}$. The only nontrivial case that we need to check is that one $D([at^m,bt^n])=[D(at^m),bt^n]+[at^m,D(bt^n)]$. But
\begin{equation}
    [D(at^m),bt^n]+[at^m,D(bt^n)]=-\langle a,b\rangle(m+n)\delta_{m,-n}K=0
\end{equation}
and $D([at^m,bt^n])=\langle a,b\rangle\delta_{m,-n-1}D(K)=0$. Since $D$ is a derivation of the Lie superalgebra $\hat{A}$, it can be extended to a derivation of its universal enveloping algebra $\mathcal{U}(\hat{A})$. Its expression is the usual one:
\begin{equation}
    D(x_{1}\cdots x_{k})=\sum_{i=1}^k x_{1}\cdots D(x_{i})\cdots x_{k},\quad x_{i}\in\hat{A}\;(1\leq i\leq k)
\end{equation}
and it's well defined because $D$ is a derivation of $\hat{A}$, i.e. $D(x\otimes y-p(x,y)y\otimes x-[x,y])=0\;\forall x,y\in\hat{A}$. \par
 We have the subalgebra $\hat{A}_+=A[t^+]\oplus\bb{C}K\subset \hat{A}$. Define the $\hat{A}_+-$module $\bb{C}C$ with the following action $at^mC=0$ for $m\geq0$ and $KC =C$ (since $K\in Z(\hat{A})$, this implies that $K$ is the identity map). Now, we induce it to a representation of $\hat{A}$ by defining the $\hat{A}-$module $F(A)=\mathcal{U}(\hat{A})\otimes_{\mathcal{U}(\hat{A}_+)}\bb{C}C$. We claim that $F(A)$ has the structure of a vertex algebra with vacuum $|0\rangle=1\otimes C$, translation operator $T=D\otimes \mathds{1}_{\bb{C}C}$ and quantum fields
 \begin{equation}
     \mathcal{F}=\left\{a(z)=\sum_{n\in\bb{Z}}(at^n)z^{-n-1}\Bigg|a\in A\right\}.
 \end{equation}
 Now, lets verify the vertex algebras axioms for $F(A)$. For the vacuum axiom, just note that $T|0\rangle = D(1)\otimes C=0$ and for $\mu\in\mathcal{U}(\hat{A})$ we have
 \begin{equation}
 \begin{split}
      [T,at^m]\mu|0\rangle &= (D(at^m\mu)-at^mD(\mu))\otimes C\\
      &=(-mat^{m-1}\mu+at^mD(\mu)-at^mD(\mu))\otimes C\\
      &=(-mat^{m-1})\mu|0\rangle.
 \end{split}
 \end{equation}
 Therefore $[T,at^m]=-mat^{m-1}\Rightarrow [T,a(z)]=\partial_za(z)$ which proves the translation covariance axiom. From the PBW-theorem applied for Lie superalgebras~(\ref{Theorem: PBW theorem for Lie superalgebras}), we have that $F(A)$ is spanned by
 \begin{equation}
     a_1t^{-m_1}\cdots a_kt^{-m_k}\otimes C\quad a_i\in A, m_i>0
 \end{equation}
 which proves the completeness axiom. Moreover, for each $a_1t^{-m_1}\cdots a_kt^{-m_k}\otimes C\in F(A)$, the commutator formula for the Clifford affinization~(\ref{Equation: commutator in the clifford afinization}) yields
 \begin{equation}
     \begin{split}
         at^n (a_1t^{-m_1}\cdots a_kt^{-m_k}|0\rangle)&= a_1t^{-m_1}\cdots a_kt^{-m_k}at^n\otimes C\\
         &=a_1t^{-m_1}\cdots a_kt^{-m_k}\otimes at^nC\\
         &=0
     \end{split}
 \end{equation}
 for $n\gg 0$ (or more precisely, for $n>\text{Max}\{m_i\}$). Therefore $\mathcal{F}$ is a collection of quantum fields. It remains to show that the quantum fields in $\mathcal{F}$ are mutually local, but
 \begin{equation}
 \begin{split}
      [a(z),b(w)]&=\sum_{m,n\in\bb{Z}}[at^m,bt^n]z^{-m-1}w^{-n-1}\\
      &=\sum_{m,n\in\bb{Z}}\langle a,b\rangle\delta_{m,-n-1}Kz^{-m-1}w^{-n-1}\\
      &=\langle a,b\rangle K\sum_{m\in\bb{Z}}z^{-m-1}w^m\\
      &=\langle a,b\rangle K\delta(z,w)
 \end{split}
 \end{equation}
 which is local.\par
 The equation equation $[a(z),b(w)]=\langle a,b\rangle K\delta(z,w)$ together with the Borcherds identity~(\ref{equation: Borcherds identity}) implies that $a_{(0)}b=\langle a,b\rangle$ (since $K$ acts on $F(A)$ as the identity map, we will omit it here). Therefore, identifying $a$ with $at^{-1}\otimes |0\rangle$, we've
 \begin{equation}
     [a_\lambda b]=\langle a,b\rangle|0\rangle.
 \end{equation}
 In expressions of this type, it will be more convenient to just write $ [a_\lambda b]=\langle a,b\rangle$, i.e. 1 for the vacuum. 
 \begin{dang}
     \begin{notation}
         It's important to emphasize that normally we will continue writing $a(z)$ for the fields $Y(at^{-1}|0\rangle,z)$. Also, in lambda-bracket relations, the vacuum $|0\rangle$ will be omitted for convenience. Maybe also note that $a_{(n)}=at^n\big|_{F(A)}$.
     \end{notation}
 \end{dang}

\subsection{The oscillator algebra}
\begin{dang}
    \begin{definition}\cite{kac2013bombay}
        The \textit{oscillator algebra} is the infinite dimensional Lie algebra $\mathcal{A}$ with basis $\{a_{m}\}_{m\in\bb{Z}}\cup\{\mathds{1}\}$ and relations $[a_m,a_n]=m\delta_{m,-n}\mathds{1}$, and $\mathds{1}$ central.
    \end{definition}
\end{dang}
There is a natural class of representations of the oscillator algebra, called the \textit{bosonic Fock space}.
\begin{dang}
    \begin{definition}\cite{kac2013bombay}
        let $m$ be a complex number and $\{\epsilon_j\;|\;j>0\}$ be a set of nonzero complex numbers. The Fock space is a representation of the oscillator algebra $\mathcal{A}$, consisting of a vector space $B^m=\bb{C}[x_1,x_2,\dots]$ of polynomials in infinitely many variables and action:
        \begin{equation}
            \begin{split}
                &\text{For }\;j >0\\
                &a_j\mapsto \epsilon_j\frac{\partial}{\partial x_j}\\
                &a_{-j}\mapsto \epsilon_j^{-1}jx_j\;\;(\text{multiplication by }\epsilon_j^{-1}jx_j)\\
                &a_0\mapsto m\mathds{1}_B\\
                &\mathds{1}\mapsto \mathds{1}_B.
            \end{split}
        \end{equation}
    \end{definition}
\end{dang}

\begin{dang}
    \begin{remark}
        It's important to note that the $\epsilon_j$ can be absorbed into redefinition of the variables $x_j$. However, we will be studying multiple actions of $\mathcal{A}$. So it's useful to keep the $\epsilon_j$ in the notation.
    \end{remark}
\end{dang}


\begin{dang}
    \begin{theorem}\cite{kac2013bombay}\label{Theorem: Fock space is irreducible}
        The representation $B^m$ of the oscillator algebra $\mathcal{A}$ is irreducible.
    \end{theorem}
\end{dang}
Proof: For any partition $\lambda=(1^{n_1}2^{n_2}\dots)$, we have a monomial $x_\lambda = x_1^{n_1}\cdots x_k^{n_k}\in B^m$ and an operator $a_1^{n_1}\cdots a_k^{n_k}$ which acts on $B^m$ as the operator $x^*_\lambda=\prod_{j=1}^k\epsilon_j^{n_j} \frac{\partial^{n_j}}{\partial x_j^{n_j}}$ on $B^m$. Define a partial order on $\mathcal{P}$ by putting $\lambda<_m\mu$ whenever $m_j(\lambda)\leq m_j(\mu)\;\forall j$.\par
It's clear that $x^*_\mu(x_\lambda)=\delta_{\mu\lambda}\prod_{j\geq1}\epsilon_j^{m_j(\mu)}m_j(\mu)!$ for any two partitions $\mu,\lambda$ with a different set of parts or satisfying $\lambda<_m\mu$. Now, suppose that $W\subset B^m$ is an invariant subspace and $v\in W$. Then we can write $v=\sum_{\lambda\in S}\alpha_\lambda x_\lambda$ where $\alpha_\lambda\in\bb{C}$ and $S\subset\mathcal{P}$ is finite. Let $\mu\in S$. Since $S$ is finite, we can admit that $\mu$ is maximal with relation to the partial order $<_m$ on $S$. Therefore
\begin{equation}
    x_\mu^*(v)=\alpha_\mu\prod_{j\geq1}\epsilon_j^{m_j(\mu)}m_j(\mu)! \in W\Rightarrow
    1\in W
\end{equation}
and by applying $a_{-j},\; j>0$ to $1$, we can reconstruct any monomial of $B^m$. This shows that $W=B^m$ and that $B^m$ is irreducible.
\begin{dang}
    \begin{theorem}\cite{kac2013bombay}\label{Theorem: Characterizaton of reps of the oscillator algebra with singular vectors}
        Let $V$ be a representation of the oscillator algebra $\mathcal{A}$. Suppose that there exists $v\in V$ such that $a_0(v)=m v,\;m\in\bb{C},\mathds{1}(v)=v$ and $a_n(v)=0,\;n>0$ (in this case $v$ is said to be a \textit{singular} vector). Then the map $B^m\rightarrow V:\prod_{j\geq1}\left(\frac{j}{\epsilon_j}\right)^{n_j} x_1^{n_1}\cdots x_k^{n_k}\mapsto a_{-1}^{n_1}\cdots a_{-k}^{n_k}(v)$ is injective. In particular, monomials of the form $a_{-1}^{n_1}\cdots a_{-k}^{n_k}(v)$ are linearly independent.
    \end{theorem}
\end{dang}
Proof: Let $W=\text{Span}\{a_{-1}^{n_1}\cdots a_{-k}^{n_k}(v)\;|n_j\geq0,k\geq0\}$. For 
 $n>0$, let's see how $a_j$ acts on a power $a_{-j}^{n_j}(v)$. Applying the commutation relation $[a_j,a_{-j}]=j\mathds{1}\;(j>0)$ repeatedly, we get
\begin{equation}
    \begin{split}
        a_ja_{-j}^{n_j}(v) &= a_{-j}a_ja_{-j}^{n_j-1}(v)+ja_{-j}^{n_j-1}(v)\\
        &=a_{-j}^2a_ja_{-j}^{n_j-2}(v)+j2a_{-j}^{n_j-1}(v)\\
        &=\cdots\\
        &=a_{-j}^{n_j}a_j(v)+jn_ja_{-j}^{n_j-1}(v)\\
        &=jn_ja_{-j}^{n_j-1}(v).
    \end{split}
\end{equation}
Now, since $[a_j,a_i]=0$ for $i+j\neq0$, we get the action $a_j=j\frac{\partial}{\partial a_{-j}}$ on $W$ for $j>0$. By the same reason, for $j<0$, $a_{-j}$ acts on $W$ just by multiplication. Therefore, we've a homomorphism of $\mathcal{A}-$modules 
\begin{equation}
    f:B^\mu\rightarrow V:\prod_{j\geq1}\left(\frac{j}{\epsilon_j}\right)^{n_j}x_1^{n_1}\cdots x_k^{n_k}\mapsto a_{-1}^{n_1}\cdots a_{-k}^{n_k}(v)
\end{equation}
 with nonzero image. Therefore, by theorem~(\ref{Theorem: Fock space is irreducible}) and Schur's lemma, $\text{ker}(f)=\{0\}$ which implies the claim of the theorem.\par
An important observation about the previous proof is that if $W=V$, then $V$ is isomorphic to $B^\mu$.\par

Now define the anti-linear map $\omega:a_n\mapsto a_{-n},\;\mathds{1}\mapsto \mathds{1}$. By anti-linear we mean $\omega(\alpha a_n)=\overline{\alpha}a_{-n}$. This map is an anti-involution, i.e. it's an involution $\omega^2=\mathds{1}_{\mathcal{A}}$ and an anti-homomorphism: 
\begin{equation}
    \omega([a_n,a_{m}])=[\omega(a_m),\omega(a_{n})]\;\forall m,n.
\end{equation}
Therefore $\omega$ can be extended to an anti-homomorphism of $\mathcal{U}(\mathcal{A})$. \begin{dang}
    \begin{notation}
       Let $V$ be a representation of the oscillator algebra with singular vector $v$. For a partition $\lambda=(\lambda_1,\dots,\lambda_r)$, it will be useful to write $a_{\pm\lambda}(v)=a_{\pm\lambda_1}\cdots a_{\pm\lambda_r}(v)$. In particular, if $\lambda =(1^{n_1},\dots ,k^{n_k})$, we have
       \begin{equation}
           a_{-\lambda}(v)=a_{-1}^{n_1}\cdots a_{-k}^{n_k}(v).
       \end{equation}
    \end{notation}
\end{dang}

\begin{dang}
    \begin{theorem}\cite{kac2013bombay}\label{Theorem: hermitian form on the representation of the oscillator algebra V}
        Let $V$ be a representation of the oscillator algebra $\mathcal{
        A}$ generated by a singular vector $v$. Then, there is a unique hermitian form $\langle\cdot|\cdot\rangle$ on $V$ satisfying the properties:
        \begin{enumerate}
            \item[(1)] $\langle\cdot|\cdot\rangle$ is contravariant with respect to $\omega$, i.e. $\langle a\psi|\psi^\prime\rangle = \langle\psi|\omega(a)\psi^\prime\rangle.$
            \item[(2)] $\langle v|v\rangle =1$.
        \end{enumerate}
        Furthermore the monomials $a_{-1}^{n_1}\cdots a_{-k}^{n_k}(v)$ are orthogonal under $\langle\cdot|\cdot\rangle$ and
        \begin{equation}\label{Equation: norm of a monomial}
            \langle a_{-1}^{n_1}\cdots a_{-k}^{n_k}(v)|a_{-1}^{n_1}\cdots a_{-k}^{n_k}(v)\rangle =\prod_{j=1}^kj^{n_j}j!.
        \end{equation}
    \end{theorem}
\end{dang}
Proof: For uniqueness, let $a_{-1}^{n_1}\cdots a_{-k}^{n_k}(v),\;a_{-1}^{m_1}\cdots a_{-k}^{m_k}(v)\;(n_i,m_i\geq0)$ be two monomials in $V$. If they aren't equal, 
we can assume that there exists an $i,\;(1\leq i\leq k)$ such that $n_i>m_i$ without lost of generality. In this case, denote $\Tilde{n}_j = n_j$ for $j\neq i$ and $\Tilde{n}_i =0$. Then, using (1), we get
\begin{equation}
    \begin{split}
        \langle a_{-1}^{n_1}\cdots a_{-k}^{n_k}(v)|a_{-1}^{m_1}\cdots a_{-k}^{m_k}(v)\rangle &= \langle a_{-i}^{n_i} a_{-1}^{\Tilde{n}_1}\cdots a_{-k}^{\Tilde{n}_k}(v)|a_{-1}^{m_1}\cdots a_{-k}^{m_k}(v)\rangle\\
         &= \langle  a_{-1}^{\Tilde{n}_1}\cdots a_{-k}^{\Tilde{n}_k}(v)|a_i^{n_i}\cdot a_{-1}^{m_1}\cdots a_{-k}^{m_k}(v)\rangle\\
          &= \langle  a_{-1}^{\Tilde{n}_1}\cdots a_{-k}^{\Tilde{n}_k}(v)|i^{n_i}\frac{\partial^{n_i}}{\partial a_{-i}^{n_i}}\cdot a_{-1}^{m_1}\cdots a_{-k}^{m_k}(v)\rangle\\
          &=0
    \end{split}
\end{equation}
because $n_i>m_i$. Moreover
\begin{equation}
\begin{split}
   a_{1}^{n_1}\cdots a_{k}^{n_k}  a_{-1}^{n_1}\cdots a_{-k}^{n_k}(v)&=\prod_{j\geq1}j^{n_j}\frac{\partial^{n_j}}{\partial a_{-j}^{n_j}}a_{-1}^{n_1}\cdots a_{-k}^{n_k}(v)\\
   &=\prod_{j\geq1}j^{n_j}n_j! v.
\end{split}
\end{equation}
 By utilizing equations (1) and (2), we obtain equation~(\ref{Equation: norm of a monomial}). To establish existence, define the inner product $\langle \cdot, \cdot \rangle$ such that the monomials $a_{-1}^{n_1} \cdots a_{-k}^{n_k}(v)$ are orthogonal, while simultaneously ensuring that equation~(\ref{Equation: norm of a monomial}) holds. So immediately we have $\langle v|v\rangle =1$. Now let $\lambda = (1^{n_1},\dots, k^{n_k}),\mu =(1^{m_1},\dots, k^{m_k})$ and $\psi=a_{-\lambda}(v),\;\psi^\prime = a_{-\mu}(v)$. If $j\geq 1$, we want to show that
 \begin{equation}\label{Equation: contravariance of the hermitian form}
     \langle a_{-j}\psi|\psi^\prime\rangle = \langle\psi|a_j\psi^\prime\rangle.
 \end{equation}
 There are two cases: If $j>k$, then $ \langle a_{-j}\psi|\psi^\prime\rangle=0$ because $a_{-j}\psi\neq \psi^\prime$. Also $a_j\psi^\prime=a_{-\mu}a_j(v)=0$, so we get equation~(\ref{Equation: contravariance of the hermitian form}). Now, if $1\leq j\leq k$, equation~(\ref{Equation: contravariance of the hermitian form}) is equivalent to the identity
 \begin{equation}
     \delta_{\Tilde{\lambda}\mu}z_\mu=jm_j\delta_{\lambda\Tilde{\mu}}z_\lambda
 \end{equation}
where $\Tilde{\lambda}=(1^{n_j},\dots,j^{n_j-1},\dots),\;\Tilde{\mu}=(1^{m_1},\dots,j^{m_j+1},\dots)$ and $z_\nu=\prod_{j\geq1}j^{m_j(\nu)}m_j(\nu)!$ for $\nu\in\mathcal{P}$. But it's clear that $\Tilde{\lambda}=\mu\iff \lambda=\Tilde{\mu}$ and in case of the equality, that $z_\mu=jm_jz_\lambda$.


 
\subsection{Charged free fermions}
Let $A=\bb{C}\psi^+\oplus\bb{C}\psi^-$ be a purely odd vector superspace of dimension 2. We put on $\mathcal{A}$ a skew-supersymmetric bilinear form $\langle\cdot,\cdot\rangle$ with matrix $\left(\begin{smallmatrix} 0&1\\1&0\end{smallmatrix}\right)$ relative to the basis $\{\psi^+,\psi^-\}$. In the last section an $\hat{A}-$module $F=F(A)$ was constructed. It's a vertex algebra with fields satisfying the relations
\begin{equation}\label{Equation: lambda-relations for charged free fermions}
     [\psi^{\pm}_\lambda\psi^\pm]=0, \quad [\psi^{\pm}_\lambda\psi^\mp]=1\;\text{ and }\; [K_\lambda -]=0.
\end{equation}
\begin{dang}
    \begin{definition}\cite{kac2013bombay}
        The \textit{bozonization} of charged free fermions is the normal ordered product $\alpha =:\psi^+\psi^-:$. 
    \end{definition}
\end{dang}
It's important to note that $p(\alpha)=0$.
\begin{dang}
    \begin{lemma}
    The bozonization of charged free fermions satisfies the identities:
    \begin{equation}
        (1)\;[\psi^\pm_\lambda\alpha]=\mp\psi^\pm\;\text{ and }\;\;(2)\;[\alpha_\lambda \alpha]=\lambda.
    \end{equation}
\end{lemma}
\end{dang}
Proof: The first identity is immediate by the noncommutative formula~(\ref{Theorem: lambda bracket identities}) and relations~(\ref{Equation: lambda-relations for charged free fermions}). Now, using the skew symmetry of the lambda bracket we get $[\alpha_\lambda\psi^\pm]=\pm\psi^\pm$. By the noncommutative Wick formula
\begin{equation}
    \begin{split}
        [\alpha_\lambda\alpha]&= :[\alpha_\lambda\psi^+]\psi^-:+p(\alpha,\psi^+):\psi^+[\alpha_\lambda\psi^-]:+\int_0^\lambda[[\alpha_\lambda\psi^+]_\mu\psi^-]d\mu\\
        &=:\psi^+\psi^-:+:\psi^+(-\psi^-):+\int_0^\lambda[\psi^+_\mu\psi^-]d\mu\\
        &=\int_0^\lambda 1d\mu= \lambda.
    \end{split}
\end{equation}
It's important to note that $[\alpha_\lambda\alpha]=\lambda$ implies $\alpha_{(1)}\alpha = |0\rangle$ and $\alpha_{(j)}\alpha=0,\;j\geq0,j\neq1$. Then $[\alpha(z),\alpha(w)]=\mathds{1}_F\partial_w\delta(z,w)\Rightarrow [\alpha_{(m)},\alpha_{(n)}]=m\delta_{m,-n}\mathds{1}_F$, i.e. $\mathcal{A}=\{\alpha_{(m)}\}_{m\in\bb{Z}}\cup\{\mathds{1}_F\}$ is the oscillator algebra.
\begin{dang}
    \begin{lemma}\label{Lemma: action of L on alpha}
        Define $L=\frac{1}{2}:\alpha\alpha:$, then the following identities are satisfied
        \begin{equation}
        (1)\;[L_\lambda\alpha]=(T+\lambda)\alpha\;\;\text{ and }\;\;(2)\; [L_\lambda L]=(T+2\lambda)L+\frac{\lambda^3}{12}.
        \end{equation}
    \end{lemma}
\end{dang}
\begin{dang}
    \begin{remark}
      We remark that equation (2) means that $L(z)$ is a Virasoro formal distribution with central charge 1.
    \end{remark}
\end{dang}
Proof: First we use the noncommutative Wick formula to compute
\begin{equation}
\begin{split}
     [\alpha_\lambda :\alpha\alpha:] &= :[\alpha_\lambda\alpha]\alpha:+p(\alpha,\alpha):\alpha[\alpha_\lambda\alpha]:+\int_0^\lambda[[\alpha_\lambda\alpha]_\mu\alpha]d\mu\\
     &=2\lambda\alpha+\int_0^\lambda[\lambda_\mu\alpha]d\mu\\
     &=2\lambda\alpha\quad(\text{because } [1_\mu\alpha]=0)\\
     &\;\Rightarrow[\alpha_\lambda L]=\lambda \alpha.
\end{split}
\end{equation}
Now, using the skew symmetry of the lambda bracket we deduce $[L_\lambda\alpha]=(T+\lambda)\alpha$. Applying again the noncommutative Wick formula and corollary~(\ref{Corollary: T acts as a derivaiton of the normal ordered product}) we get
\begin{equation}
    \begin{split}
        [L_\lambda:\alpha\alpha:]&=:[L_\lambda\alpha]\alpha:+p(L,\alpha):\alpha[L_\lambda\alpha]:+\int_0^\lambda[[L_\lambda\alpha]_\mu\alpha]d\mu\\
        &=:(T+\lambda)\alpha\alpha:+:\alpha(T+\lambda)\alpha:+\int_0^\lambda[(T+\lambda)\alpha_\mu\alpha]d\mu\\
        &=2\lambda:\alpha\alpha:+:(T\alpha)\alpha:+:\alpha (T\alpha):+\int_0^\lambda(-\mu^2+\lambda\mu )d\mu\\
        &=4\lambda L+2TL+\frac{\lambda^3}{6}
    \end{split}
\end{equation}
which proves the claim.

\begin{dang}
    \begin{lemma}\label{Equation: State field for L}
        The state field correspondence yields 
        \begin{equation}
           L= s(L(z))=\frac{1}{2}(:(T\psi^+)\psi^-:+:(T\psi^-)\psi^+:).
        \end{equation}
    \end{lemma}
\end{dang}
Proof: Using equation~(\ref{Equation: normal ordered product between quantum fields}) and the commutation relations between $\psi^+(z), \psi^-(z)$ we get
\begin{equation}
\begin{split}
    \alpha_{(-1)}&=\sum_{m<0}\psi^+_{(m)}\psi^-_{(-m-2)}+\psi^-_{(-m-1)}\psi^+_{(-m-1)}\\
    &=\sum_{m\in\bb{Z}}\psi^+_{(m)}\psi^-_{(-m-2)}.
\end{split}
\end{equation}
But, $s(L(z))=\frac{1}{2}\alpha_{(-1)}\alpha=\frac{1}{2}\alpha_{(-1)}^2|0\rangle$. Now, since $\psi^\pm_{m}|0\rangle=0$ for $m\geq0$ we have $\alpha_{(-1)}|0\rangle=\psi^-_{(-1)}\psi^+_{(-1)}|0\rangle$ and 
\begin{equation}
\begin{split}
     \alpha_{(-1)}^2|0\rangle&=\alpha_{(-1)}\psi^-_{(-1)}\psi^+_{(-1)}|0\rangle\\
    &=\psi_{(-2)}^-\psi_{(-1)}^+|0\rangle+\psi^+_{(-2)}\psi^-_{(-1)}|0\rangle\\
    &=(T\psi^-)_{(-1)}\psi^++(T\psi^+)_{(-1)}\psi^-\;\;(\text{by Lemma~(\ref{Lemma: T(anb)=Tanb+anTb and Tan=-na(n-1)})})\\
    &=:(T\psi^-)\psi^+:+:(T\psi^+)\psi^-:
\end{split}
\end{equation}
which concludes the proof.
\begin{dang}
    \begin{lemma}\label{Lemma: action of L on psi}
        The following identity is satisfied: $[L_\lambda\psi^\pm]=(T+\lambda/2)\psi^\pm$.
    \end{lemma}
\end{dang}
Proof: Using the relations~(\ref{Equation: lambda-relations for charged free fermions}) and Theorem~(\ref{Theorem: lambda bracket identities}) we deduce
\begin{equation}
    [\psi^\pm_\lambda T\psi^\pm]=0\;\;\text{ and }\;\;[\psi^\pm_\lambda T\psi^\mp]=\lambda.
\end{equation}
Now, using the noncommutative formula we get
\begin{equation}
\begin{split}
     &[\psi^+_\lambda :(T\psi^+)\psi^-:]=-T\psi^+,\;\;[\psi^+_\lambda:(T\psi^-)\psi^+:]=\lambda\psi^+,\\
      &[\psi^-_\lambda :(T\psi^+)\psi^-:]=\lambda\psi^+\;\;\text{ and }\;\;[\psi^-_\lambda:(T\psi^-)\psi^+:]=-T\psi^-.
\end{split}
\end{equation}
To get the correct equations, it's important to note that $p(\psi^\pm,\psi^\pm)=p(T\psi^\pm,\psi^\pm)=-1$. Therefore $[\psi^\pm_\lambda L]=\frac{1}{2}(\lambda-T)\psi^\pm\Rightarrow [L_\lambda\psi^\pm]=(T+\frac{1}{2}\lambda)\psi^\pm$ by the skew symmetry identity.
\begin{dang}
    \begin{definition}\cite{kac2013bombay}
        Let $(V,|0\rangle,T,\mathcal{F})$ be a vertex algebra and $L(z)\in\mathcal{F}^\prime$ a Virasoro formal distribution, i.e.,
        \begin{equation}
            [L_\lambda L]=(T+2\lambda)L+\frac{\lambda^3}{12}C.
        \end{equation}
        A field $a(z)$ is said to have of \textit{conformal weight} $\Delta\in\bb{C}$ if 
        $[L_\lambda a]=(T+\Delta\lambda)a+o(\lambda)$.
    \end{definition}
\end{dang}
Lemmas~(\ref{Lemma: action of L on alpha}) and~(\ref{Lemma: action of L on psi}) say that the the quantum fields $\psi^\pm(z), \alpha(z)$ and $L(z)$ are eigendistributions of conformal weights $1/2, 1$ and $2$ relative to $L(z)$, respectively. We will not delve into the theory of conformal weights, but just say that when a field $a(z)$ has conformal weight $\Delta$, then it's convenient to write 
\begin{equation}\label{Equation: eigendstribution of conformal weight Delta notation}
    a(z)=\sum_{n\in-\Delta+\bb{Z}}a_nz^{-n-\Delta},
\end{equation}
i.e. $a_{(n)}=a_{n-\Delta+1},\;n\in\bb{Z}$ and $a_n=a_{(n+\Delta-1)},\;n\in-\Delta+\bb{Z}$. Because then 
$[L_0,a_n]=-na_n$. Following this convention we will write
\begin{equation}\label{Equation: conformal weight convention}
    \begin{split}
        \psi^\pm(z)=\sum_{n\in-\frac{1}{2}}\psi^\pm_nz^{-n-1/2}\\
        \alpha(z)=\sum_{n\in\bb{Z}}\alpha_nz^{-n-1}\\
        L(z)=\sum_{n\in\bb{Z}}L_nz^{-n-2}.
    \end{split}
\end{equation}
Now let's summarize what we've obtained concerning these fields.
\begin{dang}
    \begin{theorem}\label{Theorem: formal distributions relations of the charged free fermions}
        Consider the Fermionic vertex algebra $F$ of charged free fermions. The following identities are satisfied:
        \begin{equation}
        \begin{split}
              &[\psi^\pm_\lambda\psi^\mp]=0,\quad[\psi^\pm_\lambda\psi^\pm]=1\\
                &[\alpha_\lambda\alpha]=\lambda,\quad [\alpha_\lambda\psi^\pm]=\pm\psi^\pm\\
                &[L_\lambda \alpha]=(T+\lambda)\alpha,\quad[L_\lambda\psi^\pm]=(T+\lambda/2)\psi^\pm.
        \end{split}
        \end{equation}
    \end{theorem}
\end{dang}
Theorem~(\ref{Theorem: formal distributions relations of the charged free fermions}) can be translated to the following commutation relations for $j>0$ 
\begin{equation}\label{Equation: nodes relations for the fermionic vertex algebra}
    \begin{split}
        &[\psi_m^\pm,\psi^\pm_n]=0,\quad[\psi^\pm_m,\psi^\mp_n]=\delta_{m+n,0}\\
        &[\alpha_m,\alpha_n]=m\delta_{m+n,0},\quad[\alpha_m,\psi^\pm_n]=\pm\psi^\pm_{m+n}\\
        &[L_m,\alpha_n]=-n\alpha_{m+n},\quad[L_m,\psi^\pm_n]=-\left(n+\frac{m}{2}\right)\psi^\pm_{m+n}.
    \end{split}
\end{equation}
In particular, we have
\begin{equation}\label{Equations: eignditribution relations}
    \begin{split}
        &(1)\;[\alpha_0,\psi^\pm_n]=\pm\psi^\pm_n\quad n\in-\frac{1}{2}+\bb{Z},\\
        &(2)\;[L_0,\alpha_n]=-n\alpha_n\quad n\in\bb{Z},\\
        &(3)\;[L_0,\psi^\pm_n]=-n\psi^\pm_n\quad n\in-\frac{1}{2}+\bb{Z}.
    \end{split}
\end{equation}

\section{The Boson-Fermion correspondence}
\subsection{The first part of the Boson-Fermion correspondence}
\begin{dang}
    \begin{theorem}\cite{kac2013bombay}
        We have a double grading $F=\bigoplus_{m\in\bb{Z}}\bigoplus_{k\geq0}F_k^m$ relative to the actions of $\alpha_0$ and $L_0$ on $F$ respectively, i.e., $\alpha_0v=mv$ and $L_ov=kv$ for $v\in F_k^m$. A state in the eigenspace $F_k^m$ is said to have \textit{charge} $m$ and \textit{energy} $k$.
    \end{theorem}
\end{dang}
Proof: By the PBW-theorem~(\ref{Theorem: PBW theorem for Lie superalgebras}) we have that $F$ is spanned by vectors of the form
\begin{equation}\label{Equation: generic monomial of F}
    \psi_{-j_t}^-\cdots\psi_{-j_1}^- \psi_{-i_s}^+\cdots\psi_{-i_1}^+|0\rangle
\end{equation}
for $j_t>\cdots>j_1>0,\;i_s>\cdots>i_1>0$ and $i_k,j_l \in -\frac{1}{2}+\bb{Z}$. By the first and third equations of~(\ref{Equations: eignditribution relations}), we immediately get
\begin{equation}
\begin{split}
      \alpha_0\psi_{-j_t}^-&\cdots\psi_{-j_1}^- \psi_{-i_s}^+\cdots\psi_{-i_1}^+|0\rangle=\\
      &=(-t+s)\psi_{-j_t}^-\cdots\psi_{-j_1}^- \psi_{-i_s}^+\cdots\psi_{-i_1}^+|0\rangle\\
        L_0\psi_{-j_t}^-\cdots&\psi_{-j_1}^- \psi_{-i_s}^+\cdots\psi_{-i_1}^+|0\rangle\\&=\left(\sum i_k+\sum j_l\right)\psi_{-j_t}^-\cdots\psi_{-j_1}^- \psi_{-i_s}^+\cdots\psi_{-i_1}^+|0\rangle.
\end{split}
\end{equation}
We write $m=s-t$ for the charge and $k=\sum i_k+\sum j_l$ for the energy.\par

\begin{dang}
    \begin{lemma}\cite{kac2013bombay}\label{Lemma: Fm has unique state of minimum energy}
        The eigenspace $F^m=\bigoplus_{k\geq0}F_k^m$ has a \textit{unique} state of minimum energy $|m\rangle$, obtained from $|0\rangle$ by applying the smallest possible number of $\psi^{\pm}$. Explicitly, one has
\begin{equation}
\begin{split}
    \left\{\begin{matrix}
        &|m\rangle =\psi^+_{-\frac{2m-1}{2}}\cdots\psi^+_{-\frac{1}{2}}|0\rangle\quad\text{if }\;m\geq0\\
        &|m\rangle =\psi^-_{-\frac{2|m|-1}{2}}\cdots\psi^-_{-\frac{1}{2}}|0\rangle\quad\text{if }\;m<0.
    \end{matrix}\right.
\end{split}
\end{equation}
    \end{lemma}
\end{dang}
Proof: Suppose that $m\geq 0$, then each monomial~(\ref{Equation: generic monomial of F})
 has energy $\sum_{k=1}^{m+t}i_k+\sum_{l=1}^tj_l$. But, both $(i_k), (j_l)$ are increasing sequences in $(-\frac{1}{2}+\bb{Z})\cap(0,\infty)$ and
 \begin{equation}
 \begin{split}
     \sum_{k=1}^{m+t}i_k+\sum_{l=1}^tj_l&\geq \sum_{k=1}^{m+t}\frac{2k-1}{2}+\sum_{l=1}^t\frac{2l-1}{2}\\
     &=\frac{(m+t)^2+t^2}{2}.
 \end{split}
 \end{equation}
 Therefore, to make energy as small as possible, it's necessary to take $t=0$ and $i_k=(2k-1)/2,\;(1\leq k\leq m)$, i.e. $|m\rangle =\psi^+_{-\frac{2m-1}{2}}\cdots\psi^+_{-\frac{1}{2}}|0\rangle$. For $m<0$, the proof is similar.
\begin{dang}
    \begin{lemma}\cite{kac2013bombay}\label{Lemma: Fm is a representation of the osillator algebra}
        $F^m$ is a representation of the oscillator algebra $\mathcal{A}=\{\alpha_n\}_{n\in\bb{Z}}\cup\{\mathds{1}_F\}$ and the state of minimum energy $|m\rangle$ is a singular vector of $F^m$.
    \end{lemma}
\end{dang}
Proof: Since $[\alpha_0,\alpha_m]=0\;\;\forall m\in\bb{Z}$, we get that $F^m$ is invariant by the action of $\alpha_n\;\forall n\in\bb{Z}$, i.e $F^m$ is an $\mathcal{A}-$module. Moreover, by the second equation in~(\ref{Equations: eignditribution relations}), $\alpha_j,\;j>0$ acts on $F^m$ by reducing the energy of its states by $j$. In particular, this proves that $\alpha_j|m\rangle=0$ for $j>0$, while $\alpha_0|m\rangle = m|m\rangle$. So $|m\rangle$ is a singular vector of the representation $F^m$ of the oscillator algebra.
\begin{dang}
    \begin{theorem}\cite{kac2013bombay}\label{Theorem: Fm is irreducible}
        For each $m\in\bb{Z}$, $F^m$ is an irreducible representation of $\mathcal{A}$.
    \end{theorem}
\end{dang}
Proof: By the lemma~(\ref{Lemma: Fm is a representation of the osillator algebra}) and theorem~(\ref{Theorem: Characterizaton of reps of the oscillator algebra with singular vectors}), there is a injective map $\eta_m:B^m\rightarrow F^m:\prod_{j\geq1}j^{n_j}x_1^{n_1}\cdots x_k^{n_k}\mapsto \alpha_{-1}^{n_1}\cdots \alpha_{-k}^{n_k}|m\rangle$ (we're taking $\epsilon_j=1$). So, because of theorem~(\ref{Theorem: Fock space is irreducible}), in order to show that $F^m$ is irreducible, it's sufficient to show that $\eta_m$ is surjective. Suppose that this is not the case, i.e. $W=F^m/\mathcal{U}(\mathcal{A})|m\rangle\neq\{0\}$. Then $W$ has a nonzero vector of minimal energy $v$ and $\alpha_j(v)=0\;\forall j>0$.\par
By the equation~(\ref{Equation: normal ordered product between quantum fields}) we get 
\begin{equation}
\begin{split}
    L_0=L_{(1)}&=\frac{1}{2}\sum_{m<0}(\alpha_m\alpha_{-m}+\alpha_{m+1}\alpha_{-m-1})\\
    &=\frac{1}{2}\alpha_0^2+\sum_{m>0}\alpha_{-m}\alpha_m\\
    &\Rightarrow L_0(v)=\frac{1}{2}m^2 v.
\end{split}
\end{equation}
So $v\in W$ is a state of energy $m^2/2$ and by lemma~(\ref{Lemma: Fm has unique state of minimum energy}) $v$ must have $|m\rangle$ as its representative, which is impossible since $|m\rangle \in \mathcal{U}(\mathcal{A})|m\rangle$ and $v\neq0$.
\begin{dang}
    \begin{notation}\label{Notation: isomorphism sigma from F to B}
        We saw in the proof of the Theorem~(\ref{Theorem: Fm is irreducible}) that $\sigma_m=\eta_m^{-1}: F^m\rightarrow B^m: \alpha_{-1}^{n_1}\cdots \alpha_{-k}^{n_k}|m\rangle \mapsto \prod_{j\geq1}j^{n_j}x_1^{n_1}\cdots x_k^{n_k} $ is an isomorphism. Therefore, we have an isomorphism
\begin{equation}
    \sigma=\bigoplus_{m\in\bb{Z}}\sigma_m: F \rightarrow B: \alpha_{-1}^{n_1}\cdots \alpha_{-k}^{n_k}|m\rangle \mapsto \prod_{j\geq1}j^{n_j}x_1^{n_1}\cdots x_k^{n_k}\mathds{1}_m
\end{equation}
where $B=\bigoplus_{m\in\bb{Z}}B^m$ and $\mathds{1}_m$ is the generator of $B^m$. To keep track of charge in $B$, it is useful to insert a formal variable $z^m$ in $B^m$ to each $m\in\bb{Z}$. So we let $B^m=z^m\bb{C}[x_1,x_2,\dots]$. In particular, this implies $\sigma_m(|m\rangle)=z^m$. 
    \end{notation}
\end{dang}


\begin{dang}
    \begin{definition}
        Let $V$ be a vector space, the \textit{formal completion} $\hat{V}$ is defined to be the vector space of $V-$valued formal distributions $V[[u^\pm]]$.
    \end{definition}
\end{dang}
\begin{dang}
    \begin{lemma}\label{Lemma: commutator of X(u) with alpha_j}
       Define $X^+(u)=u\psi^+(u), X^-(u)=\psi^-(u)$, so $X^+(u),X^-(u):F\rightarrow \hat{F}$. Then we've $[\alpha_j,X^\pm(u)]=\pm u^{ j}X^\pm(u)$. 
    \end{lemma}
\end{dang}
Proof: Using the commutation relations~(\ref{Equation: nodes relations for the fermionic vertex algebra}) we get
\begin{equation}
    \begin{split}
        [\alpha_j,X^+(u)]&=\sum_{n\in-\frac{1}{2}+\bb{Z}}[\alpha_j,\psi^+_n]u^{-n+1/2}\\
        &=\sum_{n\in-\frac{1}{2}+\bb{Z}}\psi^+_{n+j}u^{-n+1/2}\\
        &=u^jX^+(u).
    \end{split}
\end{equation}
Similarly for $X^-(u)$, $[\alpha_j,X^-(u)]=-u^{-j}X^-(u)$.\par
Now, define $\Gamma^\pm(u)=\sigma\circ X^\pm(u)\circ\sigma^{-1}$. Lemma~(\ref{Lemma: commutator of X(u) with alpha_j}) implies the commutation relations
\begin{equation}\label{Equaiton: commuation relation between Gamma+ and a_j, a_-j}
\begin{split}
    &\left[\frac{\partial}{\partial x_j},\Gamma^+(u)\right]=u^j\Gamma^+(u)\\
    &[x_j,\Gamma^+(u)]=\frac{1}{j}u^{-j}\Gamma^+(u).
\end{split}
\end{equation}
\begin{dang}
    \begin{lemma}\cite{kac2013bombay}
        We have the identities
        \begin{equation}\label{Equation: vertex operator for Gamma+}
        \begin{split}
             \Gamma^+(u)\bigl|_{\hat{B}^m}&=zu^{m+1}\exp\left(\sum_{j\geq1}x_ju^j\right)\exp\left(-\sum_{j\geq1}\frac{1}{j}\frac{\partial}{\partial x_j}u^{-j}\right)\\
              \Gamma^-(u)\bigl|_{\hat{B}^m}&=z^{-1}u^{-m}\exp\left(-\sum_{j\geq1}x_ju^j\right)\exp\left(\sum_{j\geq1}\frac{1}{j}\frac{\partial}{\partial x_j}u^{-j}\right).
        \end{split}
        \end{equation}
    \end{lemma}
\end{dang}
Proof: The reason that $z$ appears on the right side of equation~(\ref{Equation: vertex operator for Gamma+}) is because $\psi^+(u): \hat{F}^m\rightarrow \hat{F}^{m+1}$. Now, define $ T_u=\exp\left(\sum_{j\geq1}\frac{u^{-j}}{j}\frac{\partial}{\partial x_j}\right)=\prod_{j\geq1}\exp(\frac{u^{-j}}{j}\frac{\partial}{\partial x_j})$. By Lemma~(\ref{Lemma: formal taylor lemma}) $T_uf(x_1,x_2,\dots)=f(x_1+u^{-1}/1,x_2+u^{-2}/2,\dots)$, so in particular $[x_j,T_u]=-\frac{u^{-j}}{j}T_u$. Using the identity $[a,bc]=[a,b]c+b[a,c]$ we get
\begin{equation}
\begin{split}
    [x_j,\Gamma^+(u)T_u]&=[x_j,\Gamma^+(u)]T_u+\Gamma^+(u)[x_j,T_u]\\
    &=\frac{u^{-j}}{j}\Gamma^+(u)T_u+\Gamma^+(u)\left(-\frac{u^{-j}}{j}T_u\right)=0.
\end{split}
\end{equation}
Therefore, $\Gamma^+(u)T_u$ has no differential part and we can write
\begin{equation}\label{Equation: vertex 1}
    \Gamma^+(u)=zf(u,x)\exp\left(-\sum_{j\geq1}\frac{u^{-j}}{j}\frac{\partial}{\partial x_j}\right)
\end{equation}
for some unknown polynomial $f(u,x)=f(u,x_1,x_2,\dots)$. But since 
\[
\frac{\partial}{\partial x_i}\exp\left(-\sum_{j\geq1}u^jx_j\right)=-u^i\exp\left(-
\sum_{j\geq1}u^jx_j\right)
\]
we get $\left[\frac{\partial}{\partial x_i}, \exp\left(-\sum_{j\geq1}u^jx_j\right)\right]=-u^i\exp\left(-\sum_{j\geq1}u^jx_j\right)$ and
\begin{equation}
    \begin{split}
        &\left[\frac{\partial}{\partial x_i}, \exp\left(-\sum_{j\geq1}u^jx_j\right)\Gamma^+(u)\right]\\
        &\quad=\left[\frac{\partial}{\partial x_i}, \exp\left(-\sum_{j\geq1}u^jx_j\right)\right]\Gamma^+(u)+\exp\left(-\sum_{j\geq1}u^jx_j\right)\left[\frac{\partial}{\partial x_i},\Gamma^+(u)\right]\\
        &\quad=-u^i\exp\left(-\sum_{j\geq1}u^jx_j\right)\Gamma^+(u)+\exp\left(-\sum_{j\geq1}u^jx_j\right)u^i\Gamma^+(u)\\
        &\quad=0.
    \end{split}
\end{equation}
Therefore
\begin{equation}\label{Equation: vertex 2}
    \Gamma^+(u)=z\exp\left(\sum_{j\geq1}u^jx_j\right)g\left(u, \frac{\partial}{\partial x_1},\frac{\partial}{\partial x_2},\dots\right).
\end{equation}
From equations~(\ref{Equation: vertex 1}) and (\ref{Equation: vertex 2}) we get
\begin{equation}
    \Gamma^+(u)=zc_m(u)\exp\left(\sum_{j\geq1}u^jx_j\right)\exp\left(-\sum_{j\geq1}\frac{u^{-j}}{j}\frac{\partial}{\partial x_j}\right)
\end{equation}
where $c_m(u)$ still to be determined. To do this, apply 
\begin{equation}
    \Gamma^+(u)z^m=z^{m+1}c_m(u)\exp\left(\sum_{j\geq1}u^jx_j\right)\cdot 1.
\end{equation}
But 
\begin{equation}
    \begin{split}
        \Gamma^+(u)z^m&=\sigma(X^+(u)|m\rangle)\\
        &=\sum_{j\in-\frac{1}{2}+\bb{Z}}\sigma(\psi^+_j|m\rangle)u^{-j+1/2}\\
        &=\sigma(\psi_{-\frac{2(m+1)-1}{2}}|m\rangle)u^{m+1}+\cdots\\
        &=z^{m+1}u^{m+1}+\cdots
    \end{split}
\end{equation}
and so $c_m(u)=u^{m+1}$. The proof for $\Gamma^-(u)$ is totally similar.\par
Now define the maps $R(u):\hat{B}\rightarrow \hat{B}: f(x,z)\mapsto uzf(x,uz)$. Then $R(u) z^mf(x)=uz(uz)^mf(x)=zu^{m+1}\cdot z^mf(x)$. Therefore
\begin{equation}
    \Gamma^+(u) = R(u)\exp\left(\sum_{j\geq1}x_ju^j\right)\exp\left(-\sum_{j\geq1}\frac{1}{j}\frac{\partial}{\partial x_j}u^{-j}\right).
\end{equation}
Similarly, $R(u)^{-1}f(x,z)=z^{-1}f(x,zu^{-1})\Rightarrow R(u)^{-1}z^{m}f(x)=z^{m-1}u^{-m}f(x)=z^{-1}u^{-m}\cdot z^mf(x)$ and
\begin{equation}
    \Gamma^-(u)=R(u)^{-1}\exp\left(-\sum_{j\geq1}x_ju^j\right)\exp\left(\sum_{j\geq1}\frac{1}{j}\frac{\partial}{\partial x_j}u^{-j}\right).
\end{equation}
In conclusion, we have obtained the transposed action of the operators $X^\pm(u)$ to $\hat{B}$ through $\sigma$.
\begin{dang}
    \begin{theorem}\cite{kac2013bombay}
        \begin{equation}
            \begin{split}
                \Gamma^+(u) &= R(u)\exp\left(\sum_{j\geq1}x_ju^j\right)\exp\left(-\sum_{j\geq1}\frac{1}{j}\frac{\partial}{\partial x_j}u^{-j}\right)\\
                \Gamma^-(u)&=R(u)^{-1}\exp\left(-\sum_{j\geq1}x_ju^j\right)\exp\left(\sum_{j\geq1}\frac{1}{j}\frac{\partial}{\partial x_j}u^{-j}\right).
            \end{split}
        \end{equation}
    \end{theorem}
\end{dang}

\subsection{The ring of symmetric functions}
This subsection is based on~(\cite{macdonald1998symmetric}, Chapter I). The main objective of this subsection is to record some definitions that will be needed to work with symmetric functions, which play a major role in the Boson-Fermion-correspondence.\par
In the next subsection we will work with the ring of symmetric functions $\Lambda$.
Let $\Lambda_n=\bb{Z}[x_1,\dots, x_n]^{S_n}$ be the sub-ring of $\bb{Z}[x_1,\dots,x_n]$ which is invariant under the action of the symmetric group $S_n$, i.e. the set of polynomials $f$ with integer coefficients satisfying $\sigma f(x_1,\dots,x_n)=f(x_{\sigma(1)},\dots, x_{\sigma(n)})=f(x_1,\dots,x_n)$ for all $\sigma\in S_n$. \par
The ring $\Lambda_n$ is graded by degree of the polynomials. So $\Lambda_n=\bigoplus_{k\geq0}\Lambda_n^k$ 
where $\Lambda_n^k$ is the sub-group of $k$th degree symmetric polynomials. 
Let $\alpha=(\alpha_1,\dots,\alpha_n)\in\bb{N}^n$, 
we write $x^\alpha=x_1^{\alpha_1}\cdots x_n^{\alpha_n}$. If $\lambda$ is a partition with length $l(\lambda)\leq n$, 
we may complete its coordinates with zero to get an $n$-tuple. So we define
\begin{equation}
    m_\lambda(x_1,\dots,x_n)=\frac{1}{\prod_{j\geq1} m_j(\lambda)!}\sum_{\sigma\in S_n}\sigma(x^\lambda)
\end{equation}
for $l(\lambda)\leq n$, where $\sigma(x^\lambda)=x_{\sigma(1)}^{\lambda_1}\cdots x_{\sigma(n)}^{\lambda_n}$. We claim that the set $\{m_\lambda(x_1,\dots,x_n)\;|\;l(\lambda)\leq n\}$ is a basis of $\Lambda_n$. In fact, let $f=\sum_{\alpha}a_{\alpha}x^{\alpha}\in\Lambda_n$. Then
\begin{equation}
\begin{split}
     f(x_1,\dots,x_n)&=\frac{1}{n!}\sum_{\sigma\in S_n}\sigma f(x_1,\dots,x_n)\\
     &=\frac{1}{n!}\sum_{\alpha}a_\alpha\left(\sum_{\sigma\in S_n}\sigma(x^\alpha)\right).
\end{split}
\end{equation}
However, the polynomial $\sum_{\sigma\in S_n}\sigma(x^\alpha)$ is $m_\lambda(x_1,\dots,x_n)$ where $\lambda$ is obtained from $\alpha=(\alpha_1,\dots,\alpha_n)$ by listing its elements in decreasing order. Moreover, it's clear by construction that $\text{Supp}(m_\lambda(x))\cap\text{Supp}(m_\mu(x))=\varnothing$ if $\lambda\neq\mu$. Therefore, the set $\{m_\lambda(x)\;|l(\lambda)\leq n\}$ is linearly independent, and we conclude the claim.\par
From the preceding discussion, we see that $\{m_\lambda(x_1,\dots,x_n)|\;|\lambda|=k\}$ is a basis of $\Lambda_n^k$. For $m\geq n$, 
there exists a homomorphism of rings
\begin{equation}
    \rho_{mn}:\Lambda_m\rightarrow \Lambda_n: f(x_1,\dots,x_m)\mapsto f(x_1,\dots,x_n,0,\dots,0).
\end{equation}
If $l(\lambda)\leq n$, it's clear that
\begin{equation}
\rho_{mn}(m_\lambda(x_1,\dots,x_m))=m_\lambda(x_1,\dots,x_n,0,\dots,0)=m_\lambda(x_1,\dots,x_n)
\end{equation}
and if $l(\lambda)>n$, $\rho_{mn}(m_\lambda(x_1,\dots,x_m))=0$ because each monomial $\sigma(x^\lambda)$ has a factor $x_i,(n<i\leq m)$ with nonzero power. \par
Now, let $\lambda$ be a partition of $k$, then $l(\lambda)\leq k$ and the restriction $\rho_{mn}^k: \Lambda_m^k\rightarrow\Lambda_n^k$ will send basis elements to basis elements if $m\geq n\geq k$, i.e. it'll be an isomorphism. Then we define the inverse limit
\begin{equation}
\begin{split}
    \Lambda^k&=\varprojlim \Lambda_n^k\\
    &\;\;=\{(f_0,f_1,\dots,f_n,\dots)\;|\;f_n\in\Lambda_n^k\text{ and } \rho_{mn}(f_m)=f_n\text{ for }m\geq n\}.
\end{split}
\end{equation}
\begin{dang}
    \begin{example}
        We have
        \begin{equation}
            m_{(1^2)}=(0,x_1x_2, x_1x_2+x_1x_3+x_2x_3,\dots) \in \Lambda^2
        \end{equation}
        and
        \begin{equation}
            m_{(r)}=(x_1^r,x_1^r+x_2^r,x_1^r+x_2^r+x_3^r,\dots)\in\Lambda^r
        \end{equation}
        for $r\geq1$.
    \end{example}
\end{dang}
Since $\rho_{mn}$ is an isomorphism whenever $m\geq n\geq k$, the projection $\rho_n:\Lambda^k\rightarrow\Lambda_n^k$ will be an isomorphism whenever $n\geq k$. Therefore, $\Lambda^k$ is a free $\bb{Z}-$module of rank $p(k)$. Now we define the ring of symmetric functions as the direct sum $\Lambda=\bigoplus_{k\geq0}\Lambda^k$. \par
When working explicitly in the ring of symmetric functions, will take in to consideration all the partitions involved, and take a number of variables $n$, large enough to carry out all the computations in $\Lambda_n^k$.
\begin{dang}
    \begin{example}
        To multiply $m_{(1)}\cdot m_{(1)}$ in $\Lambda$, it's sufficient to work in $\Lambda_2^2$. Since $(x_1+x_2)^2=x_1^2+2x_1x_2+x_2^2$, we get 
        \begin{equation}
            m_{(1)}\cdot m_{(1)}=m_{(2)}+2m_{(1^2)}\in \Lambda.
        \end{equation}
    \end{example}
\end{dang}
The ring of symmetric functions has many interesting bases. Some of the functions that we will use in this work are the following:
\begin{align*}
    (1)&\;e_k(x_1,\dots,x_n)=\sum_{i_1<\cdots<i_k}x_{i_1}\cdots x_{i_k}\quad(\text{elementary symmetric functions}),\\
    (2)&\;h_k(x_1,\dots,x_n)=\sum_{i_1\leq\cdots\leq i_k}x_{i_1}\cdots x_{i_k}\quad(\text{complete symmetric functions}),\\
    (3)&\;p_k(x_1,\dots,x_n)=\sum_{i=1}^nx_i^k\quad(\text{power sum symmetric functions}).
\end{align*}
Elementary and symmetric functions have generating functions given by
\begin{equation}\label{Equation: generating function of elementary symmetric}
    E(u)=\sum_{i\geq0}e_iu^i= \prod_{i\geq1}(1+x_iu)
\end{equation}
and
\begin{equation}
    H(u)=\sum_{i\geq0}h_iu^i=\prod_{i\geq1}\frac{1}{1-x_iu}.
\end{equation}
For power sum symmetric functions, we have a generating function $P(u)=\sum_{i\geq1}p_iu^{i-1}$. There is a important connection between generating functions of complete and power sum symmetric functions:

\begin{equation}
    \begin{split}
        P(u)&= \sum_{i\geq 1}p_iu^{i-1}\\
        &=\sum_{i\geq1}\sum_{j\geq1}x_j^iu^{i-1}\\
        &=\sum_{j\geq1}x_j\frac{1}{1-x_ju}\\
        &=\sum_{j\geq1}\frac{d}{dt}\log\left(\frac{1}{1-x_ju}\right)\\
        &=\frac{d}{dt}\log\left(\prod_{j\geq1}\frac{1}{1-x_ju}\right)\\
        &=\frac{d}{dt}\log H(u)
    \end{split}
\end{equation}
or equivalently,
\begin{equation}\label{Equation: H(u)=exp(...)}
    H(u)=\exp\left(\sum_{j\geq1}p_ju^j/j\right).
\end{equation}

It's important to note that if we want to work with a finite number of variables, say $n$, then it's sufficient to put $x_i=0\;\forall i\geq n$. So, for example we have
\begin{equation}
    \sum_{i\geq0}\rho_n(e_i)u^i=\prod_{i=1}^n(1+x_iu).
\end{equation}
 The Schur symmetric function $s_\lambda(x_1,\dots,x_n)$ is defined by
 \begin{equation}
     s_\lambda(x_1,\dots,x_n)=\frac{\sum_{\sigma\in S_n}\text{sgn}(\sigma)\sigma(x^{\lambda+\delta})}{\prod_{i<j}(x_i-x_j)}
 \end{equation}
where $l(\lambda)\leq n$. In this work we will use 
\begin{equation}\label{Equation: jacobi-Trudy for Schur}
    s_\lambda(x_1,\dots,x_n)=\text{det}(h_{\lambda_i+j-i})_{1\leq i,j\leq n}
\end{equation}
to define the Schur symmetric functions. The equality between these two formulas is know as the Jacobi-Trudy formula.\par
The next theorem will play a central role in the final version of the Boson-Fermion correspondence (see~(\ref{Theorem: boson fermion correspondence final version})).
\begin{dang}
    \begin{theorem}\label{Theorem: the ring of symmetric functions is freely generated by the power sum}
        The ring of symmetric functions $\Lambda_{\bb{C}}=\Lambda\otimes_{\bb{Z}}\bb{C}$ is freely generated by the set of power sum and elementary symmetric functions, i.e.
        \begin{equation}
            \Lambda \cong \bb{C}[p_1,p_2,\dots]\cong\bb{C}[e_1,e_2,\dots].
        \end{equation}
    \end{theorem}
\end{dang}
Proof: See~(\cite{macdonald1998symmetric}, Chapter I).

\begin{dang}
    \begin{remark}\label{Remark: identification between the bosonic fock space and the ring of symmetric functions}
        In the next chapter we will deal with the ring of polynomials in infinitely many variables $B=\bb{C}[x_1,x_2,\dots]$. By the previous theorem, a formal change of variables $x_j\to p_j/j$ yields a identification $B\cong\Lambda_{\bb{C}}$. However, to avoid confusion we will use the symbol $S(u)=\exp(\sum_{j\geq1}x_ju^{j})=\sum_{i\geq0}S_i(x)u^i$. The polynomials $S_i(x)$ are called \textit{elementary Schur polynomials} and they are not symmetric functions. However, under the formal change $x_j\to p_j/j=(\epsilon_1^j+\epsilon_2^j+\cdots)$ we get $S_i(x_1,x_2,\dots)=h_i(\epsilon_1,\epsilon_2,\dots)$.
    \end{remark}
\end{dang}





\subsection{The second part of the Boson-Fermion correspondence}

\begin{dang}
    \begin{definition}\cite{kac2013bombay}\label{Definition: Dirac sea definition}
        Let $V=\bigoplus_{j\in\bb{Z}}\bb{C}v_j$ be a graded vector space. The \textit{Dirac sea} $\bigwedge^{\infty/2}$ is defined to be the vector space of semi-infinite monomials, i.e., formal expressions 
        \begin{equation}
            \psi =v_{i_0}\wedge v_{i_{-1}}\wedge v_{i_{-2}}\wedge\cdots
        \end{equation}
        such that the following conditions are satisfied:
        \begin{enumerate}
            \item[(1)] $i_0>i_{-1}>i_{-2}>\cdots$.
            \item[(2)] The limit $m=\lim_{k\to -\infty} (i_k-k)$ exists.
        \end{enumerate}
       The value $m$ is called the \textit{charge} of the semi-infinite monomial $\psi$.  Also, the standard rule for wedge manipulation is satisfied, i.e.
       \begin{equation}
           \cdots v_m\wedge v_n\cdots = -(\cdots v_n\wedge v_m\cdots).
       \end{equation}
    \end{definition}
\end{dang}
If a semi-infinite monomial $\psi$ has charge $m$ it will be more convenient to write
\begin{equation}
    \psi=v_{i_m}\wedge v_{i_{m-1}}\wedge\cdots
\end{equation}
so the second condition in definition~(\ref{Definition: Dirac sea definition}) is equivalent to $\lim_{k\to -\infty} (i_k-k)=0$, or $i_k=k$ for $-k\gg 0 $ which we write here as condition $(2)^\prime$.\par

Now, let $\psi=v_{i_m}\wedge v_{i_{m-1}}\wedge\cdots$, the degree of $\psi$ is defined to be the sum $\text{deg }\psi=\sum_{k=0}^\infty (i_{m-k}-m+k)$, which is a finite sum because of condition $(2)^\prime$. Moreover, for each $m\in\bb{Z}$, writing $\lambda_{k} = i_{m-k+1}+k-m-1$ implies that $\lambda = (\lambda_1,\lambda_2,\dots)$ is a partition of weight $\text{deg }\psi$. Conversely, each partition $\lambda$ of length $r$ gives a semi-infinite monomial
\begin{equation}
    \psi_\lambda^m=v_{\lambda_1-m}\wedge\cdots\wedge v_{\lambda_r-m+r-1}\wedge v_{m-r}\wedge\cdots
\end{equation}
of charge $m$ and degree $|\lambda|$. This show that partitions parametrize elements of $\bigwedge^{\infty/2}V$.
\begin{dang}
    \begin{notation}
        We will write $\psi_\lambda^m$ for semi-infinite monomials of charge $m$ and associated to the partition $\lambda$.
    \end{notation}
\end{dang}
Now, define the \textit{shift-operators} $\Lambda_kv_{j}=v_{j-k}$. For $k\neq0$, there is a well defined action of shift operators on the Dirac sea:
\begin{equation}\label{Equation: r}
     r(\Lambda_k) v_{i_1}\wedge v_{i_2}\wedge\cdots = (\Lambda_kv_{i_1})\wedge v_{i_2}\wedge\cdots + v_{i_1}\wedge(\Lambda_kv_{i_2})\wedge\cdots+\text{etc}
\end{equation}
which is a finite sum because because of the condition $(2)^\prime$. It's important to note that equation~(\ref{Equation: r}) is not well defined for $\Lambda_0$, indeed we would have
\begin{equation}
    r(\Lambda_0)v_{i_1}\wedge v_{i_2}\wedge\cdots= (1+1+\cdots)v_{i_1}\wedge v_{i_2}\wedge\cdots
\end{equation}
which doesn't make sense.
\begin{dang}
    \begin{example}
        We have
        \begin{equation}
            \begin{split}
                &\Lambda_1\psi_{(3,1)}^0=\psi_{(2,1)}^0+\psi_{(3)}^0\\
                &\Lambda_1^2\psi_{(3,1)}^0=\psi_{(1^2)}^0+2\psi_{(2)}^0\\
                &\Lambda_1^3\psi_{(3,1)}^0=3\psi_{(1)}^0\\
                &\Lambda_1^4\psi_{(3,1)}^0=3\psi_{(0)}^0.
            \end{split}
        \end{equation}
    \end{example}
\end{dang}
In general we can see that $\Lambda_k\psi_{\lambda}^m=\Lambda_1^k\psi_\lambda^m=0$ for $k\geq |\lambda|+l(\lambda)$, i.e. $\Lambda_k$ is locally nilpotent for $k>0$. Therefore, the action of $r(\Lambda_k), (k>0)$ on $\Lambda^{\infty/2}V$ can be exponentiated. We would like to define an action $R(\exp(y_k\Lambda_k))$ on $\Lambda^{\infty/2}V$ such that 
\begin{equation}\label{Equation: relation between exponention of actions on the Dirac sea}
    R(\exp y_k\Lambda_k) = \exp(r(y_k\Lambda_k))
\end{equation}
where $y_k$ is an indeterminate. To do this we will define
\begin{equation}\label{Equation: R}
    R(\exp y_k\Lambda_k)\psi = (\exp y_k\Lambda_k)v_{i_m}\wedge\cdots\wedge (\exp y_k\Lambda_k)v_{i_{m-r+1}}\wedge v_{m-r}\wedge\cdots
\end{equation}
for $\psi=v_{i_m}\wedge\cdots\wedge v_{i_{m-r+1}}\wedge v_{m-r}\wedge\cdots$ for $k>0$. It's important to note that $(\exp y_k\Lambda_k)v_j$ belongs not to $V$, but to the vector space $V[[y_k]]$. But this is not a problem since  $\Lambda_kv_j=v_{j-k}$ and $\cdots v_j\wedge\cdots\wedge v_{m-r}\cdots =0$ whenever $j\leq m-r$, so each exponential $(\cdots\wedge(\exp y_k\Lambda_k)v_j\wedge\cdots)$ inside the wedge is actually a finite sum. 
\begin{dang}
    \begin{example}\label{Example: action of exp(y1_Lambda1) on psi(3,1)}
    We have $\psi_{(3,1)}^0=v_3\wedge v_0\wedge v_{-2}\wedge v_{-3}\wedge\cdots$, and
        \begin{equation}
        \begin{split}
              R(\exp& y_1\Lambda_1)\psi_{(3,1)}^0\\
              &=\left(v_3+y_1v_2+\frac{y_1^2}{2!}v_1+\frac{y_1^3}{3!}v_0+\frac{y_1^4}{4!}v_{-1}\right)\wedge(v_0+y_1v_{-1})\wedge v_{-2}\wedge\cdots\\
              &=\psi_{(3,1)}^0+y_1(\psi_{(2,1)}^0+\psi_{(3)}^0)+\frac{y_1^2}{2!}(\psi_{(1^2)}^0+2\psi_{(2)}^0)+\frac{y_1^3}{3!}(3\psi_{(1)}^0)+\frac{y_1^4}{4!}(3\psi_{(0)}^0).
        \end{split}
        \end{equation}
    \end{example}
\end{dang}
\begin{dang}
    \begin{theorem}\cite{kac2013bombay}\label{Theorem: R(exp a)= exp(r a)}
        Let $A=\exp(y_k\Lambda_k)$ and $a=y_k\Lambda_k,\; k>0$. We have
        \begin{equation}
            R(\exp(A))=\exp(r(a)),
        \end{equation}
        where $r$ is defined by~(\ref{Equation: r}) and $R$ is defined by~(\ref{Equation: R}).
    \end{theorem}
\end{dang}
Proof: First, let $\psi = v_{i_1}\wedge v_{i_2}\wedge \cdots$ be a semi-infinite monomial with associated partition of length $r$. Then we claim that
\begin{equation}
    \frac{r(a)^k}{k!} v_{i_1}\wedge v_{i_2}\wedge \cdots = \sum_{i_1+\cdots+i_r=k}\frac{1}{i_1!\dots i_r!} a^{i_1}v_{i_1}\wedge a^{i_2}v_{i_2}\wedge \cdots.
\end{equation}
This follows by induction with base case $k=1$ being just the definition~(\ref{Equation: r}). For the induction step, suppose that the result is true for some $k\geq1$, then
\begin{equation}
    \begin{split}
        r(a)\sum_{i_1+\cdots+i_r=k}&\frac{1}{i_1!\dots i_r!} a^{i_1}v_{i_1}\wedge a^{i_2}v_{i_2}\wedge \cdots\\
        &=\sum_{i_1+\cdots+i_r=k}\frac{1}{i_1!\dots i_r!}(a^{i_1+1}v_{i_1}\wedge a^{i_2}v_{i_2}\wedge \cdots+a^{i_1}v_{i_1}\wedge a^{i_2+1}v_{i_2}\wedge \cdots+\cdots)\\
        &=\sum_{i_1+\cdots +i_r=k+1}\sum_{j=1}^r\frac{1}{i_1!\dots (i_j-1)!\dots i_r!}a^{i_1}v_{i_1}\wedge a^{i_2}v_{i_2}\wedge\cdots\\
        &=\sum_{i_1+\cdots +i_r=k+1}\frac{k+1}{i_1!\dots i_r!}a^{i_1}v_{i_1}\wedge a^{i_2}v_{i_2}\wedge\cdots
    \end{split}
\end{equation}
which proves that $\frac{r(a)^{k+1}}{(k+1)!}\psi = \sum_{i_1+\cdots +i_r=k+1}\frac{1}{i_1!\dots i_r!}a^{i_1}v_{i_1}\wedge a^{i_2}v_{i_2}\wedge\cdots$. Now
\begin{equation}
    \begin{split}
        \exp(r(a))\psi &= \sum_{k\geq0} \frac{r(a)^{k}}{k!}\psi\\
        &=\sum_{k\geq0}\sum_{i_1+\cdots +i_r=k}\frac{1}{i_1!\dots i_r!}a^{i_1}v_{i_1}\wedge a^{i_2}v_{i_2}\wedge\cdots\\
        &=\sum_{i_1,\dots,i_r}\frac{1}{i_1!\dots i_r!}a^{i_1}v_{i_1}\wedge a^{i_2}v_{i_2}\wedge\cdots\\
        &=\left(\sum_{i_1\geq0}\frac{a^{i_1}}{i_1}v_{i_1}\right)\wedge\left(\sum_{i_2\geq0}\frac{a^{i_2}}{i_2}v_{i_2}\right)\wedge\cdots \wedge \left(\sum_{i_r\geq0}\frac{a^{i_r}}{i_r}v_{i_r}\right)\wedge v_{i_{r+1}}\cdots \\
        &=R(A)\psi
    \end{split}
\end{equation}
which proves the claim.\par

Define $GL_\infty$ to be the set of infinite invertible matrices $(A_{ij})_{i,j\in\bb{Z}}$ with the property that $A_{ij}=\delta_{ij}$ whenever $|i|,|j|\geq N$ for some $N\geq1$. Multiplication is well defined in $GL_\infty$ because $(AB)_{ij}=\sum_{k=-\infty}^\infty A_{ik}B_{kj}$ is a finite sum. we also define a large group $\overline{GL}_\infty$ of infinite invertible matrices $A$ with the property that for $i\geq j$, 
there exists an integer $N\geq1$ such that $A_{ij}-\delta_{ij}=0$ whenever $|i|,|j|\geq N$. The operator $\exp y_k\Lambda_k,\; k>0$ acts on $\overline{V}$ with matrix representation in $\overline{GL}_\infty$. In fact, its matrix looks like
\begin{equation}\label{Equation: matrix representation of exp(yk_Lambdak)}
    A = \begin{pmatrix}
\ddots & \vdots & \vdots & \vdots & \vdots &\vdots & \rotatebox{30}{$\cdots$}  \\
\vdots&1 & y_k & \frac{y_k^2}{2!} & \frac{y_k^3}{3!} &  \frac{y_k^4}{3!}&\vdots \\
\vdots&0 & 1 & y_k & \frac{y_k^2}{2!} & \frac{y_k^3}{3!}&\vdots \\
\vdots& 0& 0 & 1 & y_k &\frac{y_k^2}{2!} &\vdots \\
\vdots& 0& 0 & 0 & 1 & y_k&\vdots \\
\vdots&0 & 0 & 0 & 0 & 1&\vdots \\
\rotatebox{30}{$\cdots$} & \vdots &\vdots &\vdots & \vdots & \vdots & \ddots \\
\end{pmatrix}.
\end{equation}
We can define the action of $GL_\infty$ on $\bigwedge^{\infty/2}V$ in an analogous way as for equation~(\ref{Equation: R}). But here, since $Av_{j}=v_j$ for $j\ll 0$, we just put
\begin{equation}
    R(A)v_{i_1}\wedge v_{i_2}\wedge\cdots\wedge v_{i_r}\cdots= Av_{i_1}\wedge Av_{i_2}\wedge\cdots\wedge Av_{i_r}\cdots.
\end{equation}
\begin{dang}
    \begin{theorem}\label{Theorem: R(A)psi = det A psi etc}
        Let $A\in GL_\infty$, $\psi=v_{i_m}\wedge v_{i_{m-1}}\wedge\cdots $ a semi-infinite monomial with associated partition of length $r$ and $N\geq r$  be such that $Av_{m-j}=v_{m-j}$ for $j\geq N$. Then
        \begin{equation}
        \begin{split}
            &R(A)v_{i_m}\wedge v_{i_{m-1}}\wedge \cdots\\ 
            &= \sum_{j_m>\dots>j_{m-N+1 }}\text{det}\left(A_{j_m,\dots, j_{m-N+1}}^{i_m,\dots, i_{m-N+1}}\right)v_{j_m}\wedge\cdots\wedge v_{j_{m-N+1}}\wedge v_{m-N}\cdots
        \end{split}
        \end{equation}
        where $A_{j_m,\dots ,j_{m-N+1}}^{i_m,\dots,i_{m-N+1}} = (A_{j_{m-k+1},i_{m-l+1}})_{1\leq k,l\leq N}$. 
    \end{theorem}
\end{dang}
Proof: Let $\psi=v_{i_m}\wedge\cdots \wedge v_{i_{m-r+1}}\wedge v_{m-r}\wedge\cdots$, $A\in GL_\infty$ and $N\geq r$ be such that $Av_{m-j}=v_{m-j}\;\forall j\geq N$. Then
\begin{equation}
\begin{split}
     R(A)v_{i_m}\wedge\cdots \wedge v_{i_{m-N+1}}\wedge v_{m-N}\wedge\cdots &=  Av_{i_m}\wedge\cdots \wedge Av_{i_{m-N+1}}\wedge v_{m-N}\wedge\cdots
\end{split}
\end{equation}
and
\begin{equation}\label{Equation: determinant monstrosity}
\begin{split}
     &R(A)v_{i_m}\wedge\cdots \wedge v_{i_{m-N+1}}\wedge v_{m-N}\wedge\cdots \\  &=\left(\sum_{j_m=-\infty}^\infty A_{j_mi_m}v_{j_m}\right)\wedge\cdots \wedge \left(\sum_{j_{m-N+1}=-\infty}^\infty A_{j_{m-N+1}i_{m-N+1}}v_{i_{m-N+1}}\right)\wedge v_{m-N}\wedge\cdots\\
     &=\sum_{j_m=-\infty}^\infty\cdots \sum_{j_{m-N+1}=-\infty}^\infty A_{j_mi_m}\cdots A_{j_{m-N+1}i_{m-N+1}}v_{j_m}\wedge\cdots \wedge v_{j_{m-N+1}}\wedge v_{m-N}\cdots\\
     &=\sum_{j_m>\dots >j_{m-N+1}}\left(\sum_{\sigma \in S_N}\text{sgn}\sigma A_{\sigma(j_m)i_m}\cdots  A_{\sigma(j_{m-N+1})i_{m-N+1}}\right)v_{j_m}\wedge\cdots\\
     &=\sum_{j_m>\dots >j_{m-N+1}} \left(\text{det}A_{j_m,\dots, j_{m-N+1}}^{i_m,\dots, i_{m-N+1}}\right)v_{j_m}\wedge\cdots \wedge v_{j_{m-N+1}}\wedge v_{m-N}\cdots.
\end{split}
\end{equation}



Although $\exp y_k\Lambda_k$ acts on $\overline{V}$ as a matrix in $\overline{GL}_\infty$, its action on each individual element of $\bigwedge^{\infty/2}V$ is the same as the action of some matrix in $GL_\infty$ (see example~(\ref{Example: action of exp(y1_Lambda1) on psi(3,1)})). 

\begin{dang}
    \begin{corollary}
        Let $A=\exp(y_k\Lambda_k), k>0$ and $\psi=v_{i_m}\wedge v_{i_{m-1}}\wedge\cdots$ be a semi-infinite monomial with associated partition of length $r$. Then $A$ acts on $\psi$ as an element of $GL_\infty$. Moreover, theorem~(\ref{Theorem: R(A)psi = det A psi etc}) applies to $A$ with $N=r$.
    \end{corollary}
\end{dang}
Proof: This is a direct consequence of Definition~(\ref{Equation: R}) and the fact that an exponential inside $\cdots\wedge\exp(y_k\Lambda_k)v_{i_{m-j+1}}\wedge\cdots$ is expanded at most to $y_k^r\Lambda_k^r/r!$.\par


Now, we will construct an isomorphism $F\rightarrow \bigwedge^{\infty/2}V$. To do this, for $v\in V$ and $f\in V^*$ define the \textit{wedge operator} $\Hat{v}$ and the \textit{contraction operator} $\check{f}$:
\begin{equation}
    \begin{split}
        &\hat{v}(v_{i_1}\wedge v_{i_2}\wedge\cdots) = v\wedge v_{i_1}\wedge v_{i_2}\wedge\cdots\quad\text{and}\\
        &\check{f}(v_{i_1}\wedge v_{i_2}\wedge v_{i_3}\cdots)= f(v_{i_1})v_{i_2}\wedge v_{i_3}\cdots - f(v_{i_2})v_{i_1}\wedge v_{i_3}\cdots+\cdots.
    \end{split}
\end{equation}


There is a bijection between the set of partitions $\mathcal{P}$ and the set of pair $(\alpha |\beta)$ of sequences $\alpha_1>\dots >\alpha_r\geq 0, \beta_1>\dots >\beta_r\geq 0$. The correspondence is given as follows: Given a partition $\lambda$, let $r = \text{Max}\{i\;|\lambda_i-i\geq 0\}$ be the number of squares in the main diagonal of $T_\lambda$, the diagram of $\lambda$. Now, define $\alpha_i = \lambda_i-i, \beta_i=\lambda_i^\prime -i\;(1\leq i\leq r)$ so we write $\lambda = (\alpha |\beta)$. Conversely if $(\alpha|\beta)$ is given, then $\mu=(\beta_r-r+1,\dots,\beta_1)$ is a partition and we define $\lambda=(\lambda_1,\dots,\lambda_r)$ by the equations
\begin{equation}
\begin{split}
    &\lambda_i = \alpha_i +i\quad\text{ for }\; 1\leq i\leq r \\
        &\lambda_i=\mu^\prime_{i-r}\quad \text{ for }\;r+1\leq i\leq r+\mu_1.
\end{split}
\end{equation}
With this notation, we can parametrize semi-infinite monomials directly by $\lambda=(\alpha|\beta)$ using the formula
\begin{equation}
   \psi_\lambda^m= \check{v}^*_{-\beta_r+m}\cdots\check{v}^*_{-\beta_1+m}\hat{v}_{\alpha_r+m+1}\cdots\hat{v}_{\alpha_1+m+1}\psi_0^m.
\end{equation}
Using this equality, it's possible to define a vector space isomorphism $ \bigwedge^{\infty/2}V\rightarrow F$
\begin{equation}
    \psi_\lambda^m\longmapsto \psi^-_{-\beta_r+m-\frac{1}{2}}\cdots\psi^-_{-\beta_1+m-\frac{1}{2}}\psi^+_{-\alpha_r-m-\frac{1}{2}}\cdots\psi^+_{-\alpha_1-m-\frac{1}{2}}|m\rangle.
\end{equation}
 This isomorphism is defined in such a way that the wedge operator $\hat{v}_n$ is identified with the action of $\psi^+_{-n+\frac{1}{2}}$ and the contraction operator $\check{v}_n^*$ is identified with the action of $\psi^-_{n-\frac{1}{2}}$. The reason we have introduced the Dirac sea $\bigwedge^{\infty/2}V$ is so that we can use Theorems~(\ref{Theorem: R(exp a)= exp(r a)}) and~(\ref{Theorem: R(A)psi = det A psi etc}) to do linear algebra on $\bigwedge^{\infty/2}V$.

\begin{dang}
    \begin{example}
        For charge $m=-2$ and partitions $\lambda=(1^3), (3)$ and $(2,1)$ we have the corresponding elements in $F^{-2}$ respectively: $-\psi^-_{-\frac{9}{2}}\psi^-_{-\frac{1}{2}}|0\rangle,\; \psi^-_{-\frac{5}{2}}\psi^-_{-\frac{3}{2}}\psi^-_{-\frac{1}{2}}\psi^+_{-\frac{1}{2}}|0\rangle$ and $-\psi^-_{-\frac{7}{2}}\psi^-_{-\frac{3}{2}}|0\rangle$.
    \end{example}
\end{dang}
\begin{dang}
    \begin{lemma}
        Through the isomorphism $\bigwedge^{\infty/2}V\rightarrow F$, the action of the shift 
        operator $r(\Lambda_k),\; k\neq 0$ on $\bigwedge^{\infty/2}V$ is identified with the action of $\alpha_k$ on $F$.
    \end{lemma}
\end{dang}
Proof: Using the wedge and contraction operators, it can be seen that
\begin{equation}
    r(\Lambda_k)=\sum_{j\in\bb{Z}}\hat{v}_i\check{v}_{i+k}^*
\end{equation}
which is identified with the sum $\sum_{j\in\bb{Z}}\psi^+_{-j+\frac{1}{2}}\psi^-_{j+k-\frac{1}{2}}$ acting on $F$. But using the equation~(\ref{Equation: normal ordered product between quantum fields}), the relations~(\ref{Equation: nodes relations for the fermionic vertex algebra}) and $\psi^\pm_{(n)}=\psi^\pm_{n-\frac{1}{2}}$, we get that $\alpha_k=\sum_{j\in\bb{Z}}\psi^+_{-j+\frac{1}{2}}\psi^-_{j+k-\frac{1}{2}}$.\par

As a corollary, the commutation relation $[r(\Lambda_k),r(\Lambda_l)]=k\delta_{k+l,0}$ is satisfied for $k,l\neq0$. So, although $r(\Lambda_0)$ is not defined from the definition~(\ref{Equation: r}), it's convenient to put $r(\Lambda_0)\psi_\lambda^m=m\psi_\lambda^m$, so we get a copy $\{r(\Lambda_k)\}_{k\in\bb{Z}}\cup\{\mathds{1}_F\}$ of the oscillator algebra.
\begin{dang}
    \begin{notation}
        Because of the identifications between the Dirac sea and the space $F$, we will just write $F=\bigwedge^{\infty/2}V$ with the identifications:
        \begin{align*}
            &|0\rangle \longleftrightarrow \psi_0^0\\
            &\psi^+_{-n+\frac{1}{2}} \longleftrightarrow \hat{v}_n\\
            &\psi^-_{n-\frac{1}{2}} \longleftrightarrow \check{v}^*_n\\
            &\alpha_k \longleftrightarrow r(\Lambda_k).
        \end{align*}
        We write $r^B(\Lambda_k)=\sigma\circ r(\Lambda_k)\circ\sigma^{-1}$ for the transposed action of shift operators to the bosonic Fock space $B$. Therefore (see~(\ref{Notation: isomorphism sigma from F to B}))
        \begin{equation}
            \begin{split}
                &r^B(\Lambda_k)= \frac{\partial}{\partial x_k}\;\;\text{for}\;\; k>0\\
                &r^B(\Lambda_k)=kx_k\;\;\text{for}\;\;k<0.
            \end{split}
        \end{equation}
        Also, we write $R^B(A)=\sigma\circ R(A)\circ \sigma^{-1}$ for $A\in \overline{GL}_\infty$. Consequently, by theorem~(\ref{Theorem: R(exp a)= exp(r a)}), $R^B(\exp(y_k\Lambda_k))=\exp(y_kr^B(\Lambda_k))=\exp\left(y_k\frac{\partial}{\partial x_k}\right)$.
    \end{notation}
\end{dang}
Now we can state the second part of the Boson-Fermion correspondence.
\begin{dang}
    \begin{theorem}\cite{kac2013bombay}\label{Theorem: the second part of the boson-fermion correspondence}
        Through the isomorphism
        \begin{equation}
            \sigma: F \rightarrow B : r(\Lambda_{-1})^{n_1}\cdots r(\Lambda_{-k})^{n_k}\psi_0^m\mapsto \prod_{j=1}^k j^{n_j}x_1^{n_1}\cdots x_k^{n_k}z^m
        \end{equation}
        We have
        \begin{equation}
            \sigma(v_{i_m}\wedge v_{i_{m-1}}\wedge\cdots ) = s_\lambda(x)z^m,
        \end{equation}
        where $\lambda$ is defined by $\lambda_k=i_{m-k+1}-m+k-1$ for $k\geq 1$ and $s_\lambda(x)$ is the Schur polynomial associated to the partition $\lambda$.
    \end{theorem}
\end{dang}
Proof: Since $r(\Lambda_k)$ is locally nilpotent, the operator $\prod_{k\geq 1}\exp(y_kr(\Lambda_k))$ is well defined. But since $[r(\Lambda_k),r(\Lambda_l)]=0$ for $k,l\geq 1$, $\prod_{k\geq 1}\exp(y_kr(\Lambda_k))=\exp(\sum_{k}y_kr(\Lambda_k))$. Now
\begin{equation}
    \begin{split}
        &\sigma(R(\exp(y_1\Lambda_1+y_2\Lambda_2+\cdots)v_{i_m}\wedge v_{i_{m-1}}\wedge \cdots))\\
        &\quad =R^B(\exp(y_1\Lambda_1+y_2\Lambda_2+\cdots))\sigma(v_{i_m}\wedge v_{i_{m-1}}\wedge\cdots)\\
        &\quad =R^B(\exp(y_1\Lambda_1+y_2\Lambda_2+\cdots))P(x)\\
        &\quad =\exp\left(\sum_{k\geq1}y_kr^B(\Lambda_k)\right)P(x)\\
         &\quad =\exp\left(\sum_{k\geq 1}y_k\frac{\partial}{\partial x_k}\right)P(x)\\
          &\quad =\prod_{k\geq1}\exp\left(y_k\frac{\partial}{\partial x_k}\right)P(x)\\
         &\quad =P(x+y)\quad (\text{because of lemma~(\ref{Lemma: formal taylor lemma})}).
    \end{split}
\end{equation}
Therefore $\sigma(R(\exp(y_1\Lambda_1+y_2\Lambda_2+\cdots)v_{i_m}\wedge v_{i_{m-1}}\wedge \cdots))$ is a polynomial with constant part $P(x+y)\bigl|_{x=0}=P(y)$. From the other side, this constant part must be the coefficient of $\sigma(\psi_0^m)$ in the expansion of $\sigma(R(\exp(y_1\Lambda_1+y_2\Lambda_2+\cdots)v_{i_m}\wedge v_{i_{m-1}}\wedge \cdots))$. But, by theorem~(\ref{Theorem: R(A)psi = det A psi etc}), this is $\text{det}(A_{m,m-1,\dots}^{i_m,i_{m-1},\dots})$ where $A=\exp(y_1\Lambda_1+y_2\Lambda_2+\cdots)$. Now since $\Lambda_k=\Lambda_1^k$, we have $A=\exp(y_1\Lambda_1+\cdots)=\sum_{k\geq0}S_k(y)\Lambda_k$ and $A=(A_{mn})=(h_{n-m}(y))$. So, the coefficient of $\sigma(\psi_0^m)$ is
\begin{equation}
    \begin{split}
        \text{det}(A_{m,m-1,\dots}^{i_m,i_{m-1},\dots}) &= \text{det}(A_{m-k+1,i_{m-l+1}})_{1\leq k,l\leq r}\\
        &= \text{det}(S_{i_{m-l+1}-m+k-1}(y))_{1\leq k,l\leq r}\\
        &=\text{det}(S_{\lambda_l +k-l}(y))_{1\leq k,l\leq r}\\
        &=s_\lambda(y)\quad(\text{see Remark~(\ref{Remark: identification between the bosonic fock space and the ring of symmetric functions}) and equation~(\ref{Equation: jacobi-Trudy for Schur})})
    \end{split}
\end{equation}
and the theorem is proved.
\begin{dang}
    \begin{corollary}\label{Corollary: vertex operator realization of Schur symmetric functions}
        Let 
        \begin{equation}
            S(u)= \exp\left(\sum_{j\geq1}x_ju^j\right)\exp\left(-\sum_{j\geq1}\frac{1}{j}\frac{\partial}{\partial x_j}u^{-j}\right)=\sum_{j\in \bb{Z}}S_ju^j
        \end{equation}
        and $\lambda=(\lambda_1,\dots,\lambda_r)$ be a partition. Then
        \begin{equation}
            s_\lambda(x)=S_{\lambda_1}\cdots S_{\lambda_r}\cdot 1.
        \end{equation}
    \end{corollary}
\end{dang}
Proof: It can be seen directly that $\hat{v}_k\psi_\lambda^m=\psi_{(m-k-1,\lambda)}^{m+1}$ for $k>\lambda_1+m$.  Therefore
\begin{equation}
    \hat{v}_{m+r+\lambda_1}\cdots \hat{v}_{m+1+\lambda_r}\psi_0^m=\psi_\lambda^{m+r}
\end{equation}
and by the second part of the Boson-Fermion correspondence~(\ref{Theorem: the second part of the boson-fermion correspondence}), we get
\begin{equation}
    \Gamma_{m+r+\lambda_1}\cdots \Gamma_{m+1+\lambda_r}\cdot z^m=s_\lambda(x)z^{m+r}.
\end{equation}
where $\Gamma(u)=\sum_{j\in\bb{Z}}\Gamma_ju^j$. But, by the first part of the boson-fermion correspondence, we have
\begin{equation}
\begin{split}
    \Gamma(u_r)\cdots \Gamma(u_1)\cdot z^m&=R(u_r)\cdots R(u_1)S(u_r)\cdots S(u_1)z^{m}\\
    &=u_r^{m+r}\cdots u_1^{m+1}S(u_r)\cdots S(u_1)z^{m+r}
\end{split}
\end{equation}
and comparing the coefficients of $u_{r}^{m+r+\lambda_1}\cdots u_1^{m+1+\lambda_r}$ on both sides we get 
\begin{equation}
    s_\lambda(x)z^{m+r}=\Gamma_{m+r+\lambda_1}\cdots\Gamma_{m+1+\lambda_r}\cdot z^m= S_{\lambda_1}\cdots S_{\lambda_r}\cdot z^{m+r}
\end{equation}
and $s_\lambda(x) =S_{\lambda_1}\cdots S_{\lambda_r}\cdot 1$.\par

Making a formal change of variables $x_j\rightarrow p_j/j$, where $p_j=\epsilon_1^j+\epsilon_2^j+\cdots$ denote the $j$th
power sum symmetric function, we get  the vertex realization of Schur symmetric functions:
\begin{equation}\label{Equation: vertex realization of Shcur symmetric functions}
    S(u) = \exp\left(\sum_{j\geq1}\frac{p_j}j{}u^j\right)\exp\left(-\sum_{j\geq1}\frac{\partial}{\partial p_j}u^{-j}\right).
\end{equation}
This vertex operator acts on the bosonic Fock space $B=\bb{C}[p_1,p_2,\dots]$. Observe that we can identify $B$ with the ring of symmetric functions $\Lambda_{\bb{C}}$ (see Remark~(\ref{Remark: identification between the bosonic fock space and the ring of symmetric functions}))

\begin{dang}
    \begin{theorem}\cite{jing2024characters}
    Let $S(u)$ be as in the equation~(\ref{Equation: vertex realization of Shcur symmetric functions}), then
        \begin{equation}\label{equation: Schur generating and homogeneous generatings}
            S(u_r)\cdots S(u_1)\cdot1=\prod_{1\leq i<j\leq r}\left(1-\frac{u_i}{u_j}\right)H(u_r)\cdots H(u_1).
        \end{equation}
    \end{theorem}
\end{dang}
Proof: First, we note that if $[x,y]$ commutes with $x$ and $y$, then $e^xe^y=e^ye^xe^{[x,y]}$. This follows from the identity $[x,y^k]=y^k[x,y]$ and
\begin{equation}
    \begin{split}
        xe^y &= e^yx+[x,e^y]\\
        &=e^yx+\sum_{k\geq0}\frac{1}{k!}[x,y^k]\\
        &=e^yx+\sum_{k\geq0}\frac{y^k}{k!}[x,y]\\
        &=e^y(x+[x,y]).
    \end{split}
\end{equation}
Iterating this equation, we get $x^ke^y=e^y(x+[x,y])^k\;\forall k\geq0$ and $e^xe^y=e^ye^{x+[x,y]}=e^ye^xe^{x+[x,y]}$. Now, for each $i,j\;(1\leq i<j\leq r)$, we compute
\begin{equation}
    \begin{split}
        &\exp\left(-\sum_{k\geq 1}\frac{\partial}{\partial p_k}u_j^{-k}\right)  \exp\left(\sum_{l\geq 1}\frac{1}{l}p_lu_i^{l}\right)\\
        &\quad=\exp\left(\sum_{l\geq 1}\frac{1}{l}p_lu_i^{l}\right) \exp\left(-\sum_{k\geq 1}\frac{\partial}{\partial p_k}u_j^{-k}\right)\exp\left(-\sum_{k,l\geq 1}\left[\frac{\partial}{\partial p_k},\frac{p_l}{l}\right]u_j^{-k}u_i^l\right).
    \end{split}
\end{equation}
But 
\begin{equation}
\begin{split}
      \exp\left(-\sum_{k,l\geq 1}\left[\frac{\partial}{\partial p_k},\frac{p_l}{l}\right]u_j^{-k}u_i^l\right) &=  \exp\left(-\sum_{k,l\geq 1}\frac{1}{k}\delta_{k,l}u_j^{-k}u_i^l\right)\\
      &=\exp\left(-\sum_{k\geq 1}\frac{1}{k}\left(\frac{u_i}{u_j}\right)^k\right)\\
      &=\left(1-\frac{u_i}{u_j}\right)
\end{split}
\end{equation}
while
\begin{equation}
    \exp\left(-\sum_{k\geq 1}\frac{\partial}{\partial p_k}u_j^{-k}\right)\cdot1=1
\end{equation}
for all $j$. Furthermore, by the equation~(\ref{Equation: H(u)=exp(...)})
\begin{equation}
   \prod_{i=1}^r \exp\left(\sum_{k\geq 1}\frac{1}{k}p_ku_i^{k}\right)\cdot 1=H(u_r)\cdots H(u_1).
\end{equation}
So, moving all the differential operators to the right and doing the necessary commutations we get the result.

\begin{dang}
    \begin{corollary}   
    By extracting coefficients $u^{\lambda},\;\lambda=(\lambda_1,\dots,\lambda_r)$ on both sides of 
    equation~(\ref{equation: Schur generating and homogeneous generatings})
    we get the identity
    \begin{equation}
        S_\lambda\cdot 1 = \prod_{i<j}(1-R_{ij})h_\lambda
    \end{equation}
    where $R_{ij}h_\lambda = h_{(\lambda_1,\dots,\lambda_i+1,\dots,\lambda_j-1,\dots,\lambda_r)}$. 
    \end{corollary}
\end{dang}


\subsection{Orthogonality}
There is an Hermitian form on $F$, relative to which the semi-infinite monomials are orthonormal. 
\begin{dang}
    \begin{lemma}\cite{kac2013bombay}\label{Lemma: E_ji is the adjoint of E_ij}
    Let $\langle\cdot|\cdot\rangle$ be the hermitian form on $F$ where the set of semi-infinite monomials forms an orthonormal set. For two semi-infinite monomials $\psi,\psi^\prime\in F$ we have
        \begin{equation}\label{Equation: E_ji is the adjoint of E_ij}
            \langle r(E_{ij})\psi|\psi^\prime \rangle=\langle \psi|r(E_{ji})\psi^\prime \rangle
        \end{equation}
        where $E_{ij}v_k = \delta_{jk}v_i$.
    \end{lemma}
\end{dang}
Proof: Let $\psi = v_{i_1}\wedge v_{i_2}\wedge \cdots$, there are two cases: If $j\notin\{i_1,i_2,\dots\}$, then $r(E_{ij})\psi=0\Rightarrow \langle r(E_{ij})\psi|\psi^\prime \rangle=0$. But in this case we necessarily have $\psi \neq r(E_{ji})\psi^\prime$, so $\langle \psi|r(E_{ji})\psi^\prime \rangle=0$. Now, if $j\in \{i_1,i_2,\dots\}$ then $r(E_{ij})\psi=v_{i_1}\wedge\cdots \wedge v_{j}\cdots \neq 0$ and $r(E_{ji})r(E_{ij})\psi = \psi$. Therefore $\psi^\prime = r(E_{ij})\psi \iff r(E_{ji})\psi^\prime = \psi$, and equation~(\ref{Equation: E_ji is the adjoint of E_ij}) is satisfied.\par

Since $\Lambda_k=\sum_{j\in\bb{Z}}E_{j,j+k}$, lemma~(\ref{Lemma: E_ji is the adjoint of E_ij}) says that
\begin{equation}\label{Equation: hermitian form on F is contravariant with respect to omega}
     \langle r(\Lambda_k)\psi|\psi^\prime \rangle=\langle \psi|r(\Lambda_k)^\dagger\psi^\prime \rangle
\end{equation}
where $A^\dagger$ means the transpose conjugate of $A\in \overline{GL}_\infty$. But since $r(\Lambda_k)^\dagger =r(\Lambda_{-k})=\omega(r(\Lambda_{k}))$, we get by Theorem~\ref{Theorem: hermitian form on the representation of the oscillator algebra V} that $\langle\cdot|\cdot\rangle$ is the unique hermitian form on $F^m$ which is contravariant with respect to $\omega$ satisfying $\langle \psi_0^m|\psi_0^m\rangle=1$ (remember that $r(\Lambda_k)\psi_0^m=0$ for $k>0$, i.e., $\psi_0^m$ is a singular vector of $F$).\par
Using the Boson-Fermion correspondence, we transport the Hermitian form $\langle\cdot|\cdot\rangle$ on $F^m$ to a Hermitian form $\langle\cdot|\cdot\rangle_{B^m}$ on $B^m$ through the isomorphism $\sigma_m: F^m\rightarrow B^m$. By equation~(\ref{Equation: hermitian form on F is contravariant with respect to omega}), the transported hermitian form satisfies (for $k>0$)
\begin{equation}\label{Equation: transposed hermitian form on B}
\begin{split}
     \langle r^{B^m}(\Lambda_{-k})\cdot P|Q\rangle_{B^m}&=\langle P|r^B(\Lambda_k)\cdot Q\rangle_{B^m}\iff \\
          \langle kx_k\cdot P|Q\rangle_{B^m}&=\biggl< P\bigg|\frac{\partial}{\partial x_k} Q\biggr>_{B^m}
\end{split}
\end{equation}
and $\langle1|1\rangle_{B^m}=1$. Using equation~(\ref{Equation: transposed hermitian form on B}), we get that $ \langle 1|P\rangle_{B^m}=P(x)|_{x=0}\langle 1|1\rangle_{B^m}=P(x)|_{x=0}$ and
\begin{equation}\label{Equation: transposed hermitian form to B}
\begin{split}
     \langle P|Q\rangle_{B^m}&=\biggl< 1\bigg| P\left(\frac{\partial}{\partial x_1},\frac{1}{2}\frac{\partial}{\partial x_2},\dots\right)Q\biggr>_{B^m}\\
     &= P\left(\frac{\partial}{\partial x_1},\frac{1}{2}\frac{\partial}{\partial x_2},\dots\right)Q(x)\bigg|_{x=0}.
\end{split}
\end{equation}
Under the change of variables $x_j\to p_j/j$, equation~(\ref{Equation: transposed hermitian form to B}) implies 
\begin{equation}\label{Equation: orthogonality of power sums}
    \langle p_\lambda|p_\mu\rangle_{B^m} = \delta_{\lambda\mu}z_\lambda
\end{equation}
where $p_\lambda=p_{\lambda_1}p_{\lambda_2}\cdots$ and $z_\lambda=\prod_{j\geq1}j^{m_j(\lambda)}m_j(\lambda)!$. Equation~(\ref{Equation: orthogonality of power sums}) defines a Hermitian form on the Fock space ${B^m}=\bb{C}[p_1,p_2,\dots]$. It's important to note that in this case, the action of the oscillator algebra $\mathcal{A}=\{a_n\}_{n\in\bb{Z}}\cup\{\mathds{1}\}$ is (for $n>0$):
\begin{equation}
    \begin{split}
        &a_{-n} \mapsto p_n\\
        &a_n\mapsto n\frac{\partial}{\partial p_n}\\
        &a_0\mapsto m\mathds{1}_{B^m}\\
        &\mathds{1}\mapsto \mathds{1}_{B^m}.
    \end{split}
\end{equation}
Therefore, the anti-linear anti-involution $\omega$ of $\mathcal{A}$ defines an adjoint operator $p_n^\dagger = n\frac{\partial}{\partial p_n}$. By the second part of the Boson-Fermion correspondence~(\ref{Theorem: the second part of the boson-fermion correspondence}), equation~(\ref{Equation: orthogonality of power sums}) defines the unique hermitian form with the property
\begin{equation}\label{Equation: orthonormality of Schur symmetric functions}
   \langle s_\lambda|s_\mu\rangle_{B^m} = \delta_{\lambda\mu}.
\end{equation}
The ring of symmetric functions $\Lambda_{\bb
C}$ can be identified with the Fock space $\bb{C}[p_1,p_2,\dots]$. So, equation~(\ref{Equation: orthogonality of power sums}) defines a hermitian form $\langle\cdot,\cdot\rangle$ on $\Lambda_{\bb{C}}$ and by equation~(\ref{Equation: orthonormality of Schur symmetric functions}), the set of Schur symmetric functions forms an orthonormal basis of $(\Lambda_{\bb{C}}, \langle\cdot,\cdot\rangle)$.\par
Using Corollary~(\ref{Corollary: vertex operator realization of Schur symmetric functions}) and equation~(\ref{Equation: vertex realization of Shcur symmetric functions}), get the a version of the Boson-Fermion correspondence that will be of interest in the next chapter.
\begin{dang}
    \begin{theorem}\label{Theorem: boson fermion correspondence final version}
        Let $p_j = \sum_{i\geq1}\epsilon_i^j$ be power sum functions in $\epsilon$ and $\Lambda_{\bb{C}}=\bb{C}[p_1,p_2,\dots]$ be the ring of symmetric functions. The vertex operator 
        \begin{equation}
              S(u) = \exp\left(\sum_{j\geq1}\frac{p_j}j{}u^j\right)\exp\left(-\sum_{j\geq1}\frac{\partial}{\partial p_j}u^{-j}\right)
        \end{equation}
        gives a realization of Schur symmetric functions through the formula:
        \begin{equation}
            s_\lambda(\epsilon) = S_{\lambda_1}\cdots S_{\lambda_r}\cdot 1.
        \end{equation}
        Also, the set of Schur symmetric functions form an orthonormal basis of $\Lambda_{\bb{C}}$ with respect to the hermitian form defined by
        \begin{equation}
            \langle p_\lambda,p_\mu\rangle = \delta_{\lambda\mu}z_\lambda
        \end{equation}
        where $z_\lambda = \prod_{j\geq1}j^{m_j(\lambda)}m_j(\lambda)!$. Moreover, this forms is the unique Hermitian form which is contravariant with respect to the adjoint operator $p_n^\dagger = n\frac{\partial}{\partial p_n}$.
    \end{theorem}
\end{dang}
